% \documentclass{article}
  % \usepackage{maa-monthly}
 \documentclass{amsart}
% \documentclass{siamltex}


  \let\oldchapter=\chapter
\def\resetpage{\setcounter{page}{1}}
\def\chapter{\expandafter\resetpage\oldchapter}  
 
\usepackage{mathdots}

\usepackage[all]{xy}

\usepackage{bm}
\usepackage{bigints}
\usepackage{epsfig}
\usepackage{epstopdf}

\usepackage{srcltx,amssymb}

\usepackage{graphicx}

    \usepackage{amsthm}
    \raggedbottom
      \usepackage{amsmath}
    %% IF YOU HAVE FONTS INSTALLED
    %\usepackage{mtpro2}
    \usepackage{amscd}
    %\usepackage{mathtime}\bullet
    \usepackage[title]{appendix}
    \usepackage{slashed}
    \numberwithin{equation}{section} 
       
   %\theoremstyle{theorem}
   \newtheorem{theorem}{Theorem}
    \newtheorem*{thm2}{Proposition A}
    \newtheorem*{thm5}{Proposition A$'$}
    \newtheorem*{thm3}{Proposition B}
    \newtheorem*{thm4}{Proposition B$'$}
  %  \newtheorem{corollary}[theorem]{Corollary}

\newtheorem{remark}[theorem]{Remark}

\newtheorem*{conjecture}{Conjecture}

   \usepackage{appendix}
   
 \newtheorem{lemma}[theorem]{Lemma}
  \newtheorem{corollary}[theorem]{Corollary}
  %  \newtheorem{theorem}{Theorem}
     \newtheorem{proposition}[theorem]{Proposition}
     
    \theoremstyle{definition}
   \newtheorem{definition}[theorem]{Definition}
  
    
    
   \numberwithin{theorem}{section} 
    
  %  \newtheorem{definition}[theorem]{Definition}
  
  \setcounter{tocdepth}{1}
    
    \begin{document}
       
             \title{Novel isomorphism invariants of real algebras}
    \author{Fred Greensite}
    
   \begin{abstract}  For real finite-dimensional unital associative algebras we identify a group of characterizations of an algebra's multiplicative inversion operation having features that include novel algebra isomorphism invariants, and which also forward a particular generalization of the usual algebra norm.  These provide further insights into comparative algebra structure which could aid in fashioning useful alternative systems of algebra classification that are coarser though more tractable than that based solely on isomorphism equivalence.  This unital algebra program can be augmented by addressing other unary operations dependent on the algebra product, making portions of the program applicable to nonunital algebras, and otherwise potentially enhancing the classification power of the approach.     
   \end{abstract}
        
  \maketitle
 % \markright{Novel isomorphism invariants of real algebras} 
 \markboth{F. Greensite }{Novel isomorphism invariants of real algebras}
  
 \tableofcontents
 
  \section{Introduction} The general goal of classifying algebras with respect to isomorphism equivalence is said to be hopeless according to its recognition as a so-called wild problem, e.g., \cite{drozd:1972,prest:1991,sergeichuk:2000,belitskii:2005}.   Thus, the question of identifying a less ambitious but transparent, manageable, and meaningful classification system is germane.  As inspiration, one can consider Topology's use of homotopy, homology, etc., to fashion classification systems coarser than that based on homeomorphism equivalence.  Toward this end, we forward a program based on  ``inversion characterization groups" relevant to real finite-dimensional unital algebras.  These groups express redundancies in characterization of the multiplicative inversion operation, a distinctive nonlinear unary operation implied (for example) by unital associative algebras. One application is that, together with an associated functor, this augments the means for addressing the question of whether two algebras are related by an epimorphism, via exploitation of novel computable isomorphism invariants embedded in the group structure. As for classification per se, one way this system (very coarsely) organizes algebras is according to how they are associated with invariant weighted combinations of three elemental characteristics pertaining to the above groups, which can be referred to as rational, logarithmic, and angular.  
A further feature of this program is that it generalizes the notion of an algebra norm.  
  
The possibility also exists that the forwarded classification program can be significantly enhanced by inclusion of structures distinct from (though related in concept to) those leading to inversion characterization groups.  This depends on the degree to which an algebra's ``uncurling metric groups of various types" as specified in Definition \ref{7/21/23.1} differ from each other.  Investigation of this possibility is left for the future.
  
         It is most efficient to immediately present the mechanism of our program, and that is what we will do.  However, the mechanism did not come out of nowhere, and two distinct motivations are presented as Appendix A (where the mechanism emerges from a new proof of the Pythagorean Theorem based on exploitation of an ``isotropic" feature of Euclidean space) and Appendix B (where the mechanism emerges from a geometrical interpretation of an unusual Inverse Problems regularization method).
         
   Appendix C presents an organizational chart of the relationships between topics introduced and/or addressed in this paper.     

\bigskip
 
 {\it Conventions:}   Unless stated otherwise, ``algebra" means a real finite-dimensional unital associative algebra, i.e., a unital associative algebra whose vector space of elements is $\mathbb{R}^n$. The standard topology on $\mathbb{R}^n$ is assumed, so that for analytic purposes an algebra's vector space of elements is assumed to incorporate the Euclidean norm/inner product.  For expressions involving the usual dot product ``\,$\cdot$\,", the operands are understood to be taken with respect to the standard basis. The juxtaposition of two algebra elements indicates their algebra product.   An algebra will typically be denoted by $A$, sometimes subscripted, etc. The multiplicative identity of algebra $A$ is denoted as $\mathbf{1}_A$ or just $\mathbf{1}$ if the algebra in question is understood.   As a vector, an element $s\in A$ has its dual vector denoted by $s^T$, and we will also use the notation $s^{-T} \equiv (s^{-1})^T = (s^T)^{-1}$.  The transpose of a linear operator $K$ is denoted by $K^T$.                
       
       \bigskip
       


    
 
 \section{Groups of uncurling metrics, unital norms, and associated  invariants}
 
 
   
   \subsection{The uncurling metric group of an algebra}
     
  \begin{definition}\label{7/21/23.1} With respect to elements $s$ of a given algebra $A$, where these elements are in a neighborhood of the algebra's multiplicative identity that consists only of units, a  {\it type j uncurling metric $L$}, with $j\in \mathbb{Z}$, is a symmetric linear transformation  satisfying \begin{equation}\label{4/7/23.1} d([Ls^j]\cdot ds) = d\left((s^j)^T L ds\right) = 0. \end{equation}  The additive group of all such linear transformations is the algebra's {\it uncurling metric group of type j}.   \end{definition}
  
 
 The curl is the Hodge dual of the exterior derivative of the differential \mbox{1-form} dual to a vector field (i.e., the curl is the dual of the infinitesimal rotation tensor).  Thus, (\ref{4/7/23.1}) states that the curl of $Ls^j$ vanishes ($L$ ``uncurls" $s^j$).    In other words (recalling the conventions stated in the Introduction), $L$ is an uncurling metric of type $j$ if, with respect to the standard basis, $Ls^j$ is a gradient.   Nothing obtuse results from specification of the standard basis in Definition \ref{7/21/23.1}, because uncurling metric groups associated with isomorphic algebras are intimately related - as seen in Theorem \ref{6/6/23.1}/Remark \ref{9/18/23.1}.
  Furthermore, although it is matrix addition that figures in any one of these groups in isolation, Theorem \ref{6/6/23.1} demonstrates the important role of the matrix product when distinct groups of uncurling metrics are being compared.  
   
\begin{remark}\label{10/16/23.1} {\rm  Most of this paper will concern the uncurling metric group of \mbox{type -1}, because that group leads to a generalization of the usual algebra norm (Section 3), as well as augmentation of the set of easily recognizable computationally accessible isomorphism invariants beyond those found in \mbox{Corollary \ref{8/18/23.2}} (i.e., there are the additional invariants identified in Theorem \ref{6/6/23.3}). Two more reasons are found in Appendix A and Appendix B. Of course, uncurling metric groups of negative type do not exist for the case of nonunital algebras.  The nonunital algebra setting will be briefly addressed in Section 6.  } \end{remark}
 
   
Because it figures so prominently in this paper, the uncurling metric group of type -1 will be denoted simply by $\mathfrak{u}_A$, or by $\mathfrak{u}$ if the algebra in question is understood.  Furthermore, when we refer to the ``uncurling metric group" without a type specified we mean the uncurling metric group of type -1.  Likewise, an ``uncurling metric" without a type specified refers to a member of the uncurling metric group of \mbox{type -1}. 
 
      Given the particular importance of $\mathfrak{u}_A$, a question that immediately arises is whether or not this group is always nontrivial (i.e., has more than just the matrix of zeros as a member).  This will be answered in Section 3, wherein it is shown that the uncurling metric group of type -1 is always at least one-dimensional.  Surprisingly, a further question of importance is whether all potentially interesting types of uncurling metric groups (i.e., all types with $j\ne 0,1$) are identical for a given algebra (see the first paragraph of Section 5.2).  The latter question is not answered in this paper.

\begin{remark}\label{11/21/23.1} {\rm The uncurling metric group of type 1 does not address the given algebra, since the mapping defined by $s\mapsto s$ does not depend on the algebra product.  Hence, a type 1 group would seem to be pointless given the desire to derive an entity sensitive to the algebra and not just the underlying vector space.  However, the type 1 group does forward equations (i.e., (\ref{11/29/23.1}), (\ref{11/29/23.2}), and ultimately, (\ref{11/29/23.5}), (\ref{11/29/23.6})) that can serve as the basis for a new proof of the Pythagorean Theorem, as presented in Appendix A.  It is also seen in Appendix A that, viewed in isolation, the mechanism underlying the new proof can serve as a motivation for our whole program.}\end{remark}

\enlargethispage{1cm}

The novel invariants associated with uncurling metric groups all proceed from the following.


\begin{theorem} \label{6/6/23.1} If $K:A_1\rightarrow A_2$ is an isomorphism of algebras, then \begin{equation}\label{5/23/23.1} \mathfrak{u}_{A_1} = K^T \mathfrak{u}_{A_2} K. \end{equation} \end{theorem}

\begin{proof} Let $s_1,s_2$ denote members of the spaces of units of $A_1,A_2$, respectively, with $L_2\in \mathfrak{u}_{A_2}$.  Then, \begin{equation} 0 = d\left(s_2^{-T}L_2 ds_2\right) =  d\left(Ks_1)^{-T}L_2 d(Ks_1)\right) =d\left(s_1^{-T} \big[K^T L_2K\big]ds_1\right), \label{9/21/23.1}  \end{equation} where we have used the commutativity of a differential with a linear transformation, and the fact that an algebra isomorphism maps a multiplicative identity to a multiplicative identity (i.e., ${\bf 1}_{A_2} = K\left(\mathbf{1}_{A_1}\right) = K(s_1s_1^{-1}) = K(s_1)K(s_1^{-1})$, so that $(Ks_1)^{-1} = K(s_1^{-1})$).  It follows from (\ref{9/21/23.1}) that every uncurling metric of algebra $A_2$ is associated with an uncurling metric of $A_1$, i.e., from the last equality in (\ref{9/21/23.1}) we have  $K^TL_2K\in\mathfrak{u}_{A_1}$.  An analogous argument employing $K^{-1}$ indicates that for $L_1\in\mathfrak{u}_{A_1}$ we have the associated $(K^{-1})^T L_1 K^{-1} \in \mathfrak{u}_{A_2}$.  Furthermore, $K^T L_2K\in\mathfrak{u}_{A_1}$ is such that $(K^{-1})^T [K^T L_2 K]K^{-1} = (K^T)^{-1} [K^T L_2 K]K^{-1} = L_2$.  Similarly, $(K^{-1})^T L_1K^{-1}\in\mathfrak{u}_{A_2}$ is such that $K^T[(K^{-1})^T L_1K^{-1}] K = L_1$.  The theorem follows from the implied bijection. 
\end{proof}

\begin{remark} \label{9/18/23.1} {\rm Theorem \ref{6/6/23.1} easily generalizes to uncurling metric groups of arbitrary type, based on the isomorphism condition $K(s_1^j) = (K(s_1))^j$.   } \end{remark}

Computation of $\mathfrak{u}$ is straightforward in principle.  That is, for an $n$-dimensional algebra with elements $s = (x_1,x_2,\dots,x_n)$ with respect to the standard basis, one initially takes $\mathfrak{u}$ to be a symmetric $(n\times n)$-matrix  whose upper triangular entries are each a different member of a set of $\frac{n(n+1)}{2}$ real parameters.  Computation of $s^{-1}$ produces a vector whose components depend on $x_1,x_2,\dots,x_n$.  One then writes down the components of $\mathfrak{u}s^{-1}$, i.e., an $n$-vector whose components are dependent on the parameters in the matrix $\mathfrak{u}$ and $x_1,x_2,\dots,x_n$.  There are then $\frac{n(n-1)}{2}$ constraint equations that must be satisfied in order that (\ref{4/7/23.1}) hold.  Satisfaction of each equation ultimately reduces to the requirement that a particular multivariate polynomial be {\it identically} zero.  After simplification of the polynomial to be a linear combination of independent monomials, each monomial term is then comprised of a coefficient consisting of a linear combination of some of the above $\frac{n(n+1)}{2}$ parameters, multiplying an expression consisting of a product of some of the components of $s$, each raised to some power.  Therefore, each of the above coefficients of these (independent) monomials are required to be zero.  Accordingly, the number of surviving independent parameters (i.e., parameters not required to be equated to some linear combination of other parameters, or set to 0) is successively reduced as the $\frac{n(n-1)}{2}$ constraint equations are successively treated.  The final number of surviving independent parameters is the dimension of $\mathfrak{u}$.

From this standpoint, the members of $\mathfrak{u}$ are realizations of the ``matrix of parameters" $M_{\mathfrak{u}}$, whose entries are particular linear combinations of the above surviving independent parameters (i.e., each realization/uncurling metric results from a single set of selected real values for those parameters).  The determinant of $M_{\mathfrak{u}}$, written as ${\rm det}(M_{\mathfrak{u}})$, is a polynomial in the surviving parameters.  

\enlargethispage{1cm}

   With the above in mind, it is seen that (\ref{5/23/23.1}) has immediate consequences.
    
  
    
  \begin{corollary}\label{8/18/23.2} The following are algebra isomorphism invariants. \begin{itemize} \item The dimension of $\mathfrak{u}$. \item The dimension of the subspace of the algebra's vector space of elements to which ${\rm det}(M_{\mathfrak{u}})$ is sensitive. \item The real dimension of the real algebraic set specified by ${\rm det}(M_{\mathfrak{u}}) = 0$.  \item The dimension of the smallest kernel of the members of $\mathfrak{u}$.  \item The dimension of the largest kernel of the members of $\mathfrak{u}$ exclusive of the member that maps identically to zero.\end{itemize} \end{corollary}

 It is shown in Section 3 that $\mathfrak{u}$ is always at least one-dimensional.   Concrete examples of computation of $\mathfrak{u}$ are found in Section 2.3 and Section 4.  
 
 Of course, the above corollary generalizes to uncurling metric groups of arbitrary type, and is of potential interest for types $j\ne 0,1$.

\subsection{The unital norm group of an algebra}
  
Our later Definition \ref{9/20/23.1} introduces the notion of a ``normalized" uncurling metric.  Its specification will employ the (standard basis) computation of a positive integer $\|\mathbf{1}\|^2 \equiv \mathbf{1} \cdot \mathbf{1}$. As preparation, this factor must first figure in Definition \ref{9/21/23.2} below.   Note that $\|\mathbf{1}\|^2$ is not an algebra isomorphism invariant.  That is, for isomorphic algebras $A_1,A_2$, it may be the case that $\|\mathbf{1}_{A_1}\|^2 \ne \|\mathbf{1}_{A_2}\|^2$ (in which case the isomorphism is not isometric).  This basis-specific factor does not figure in the main result of this subsection, \mbox{Theorem \ref{6/6/23.3}.}

  
  \begin{definition}\label{9/21/23.2} Let $L$ be an uncurling metric for some algebra.  For $s$ in a convex neighborhood of $\mathbf{1}$ devoid of non-units, the associated {\it unital norm} $ \ell(s)$ is the function
  \begin{equation}\label{11/7/20.2}   \ell(s) \equiv \exp\left({\frac{1}{\|{\bf 1}\|^2}\int_{\bf 1}^s [Lt^{-1}]\cdot dt}\right).\end{equation} 
   The set of unital norms associated with an algebra, which is an abelian group $U$ with respect to the usual point-wise operation of multiplication of functions, is called the {\it unital norm group}. \end{definition} 
   
To see that a convex neighborhood of $\mathbf{1}$ devoid of units exists for any algebra, it is sufficient to observe that such a neighborhood clearly exists for the left regular representation of the algebra (where one works within the matrix algebra $ M_n(\mathbb{R})$, and the inverse of an element is simply the reciprocal of the determinant multiplied by the adjugate matrix).
  
    From (\ref{11/7/20.2}), an uncurling metric and its associated unital norm satisfy  \begin{equation}\label{11/7/20.1} Ls^{-1}  = \nabla\left[\|{\bf 1}\|^2\log  \ell(s)\right] =   \left(\frac{\|{\bf 1}\|^2}{ \ell(s)}\right)\nabla \ell(s).\end{equation} 
Thus, the pair $(L, \ell(s))$ characterizes the algebra's multiplicative inversion operation $s\mapsto s^{-1}$ in the portion of a simply connected neighborhood of $\mathbf{1}$ devoid of non-units that is orthogonal to the kernel of $L$.  If $L$ is nonsingular, this pair fully characterizes $s^{-1}$ in a simply connected neighborhood of $\mathbf{1}$ devoid of non-units.  For a given algebra, the set of all such pairs (where $L$ may or may not be singular) is an abelian group with respect to the binary operation 
\begin{equation}\label{5/23/23.2} (L_1, \ell_1(s))(L_2, \ell_2(s)) \equiv 
\left(L_1+L_2, \ell_1(s)  \ell_2(s)\right).\end{equation} 
\begin{definition} For a given algebra, the set of pairs $(L, \ell(s))$ satisfying (\ref{4/7/23.1}) and (\ref{11/7/20.2}), in the context of the operation defined by (\ref{5/23/23.2}), is the algebra's {\it inversion characterization group}. \end{definition}

Any member of the inversion characterization group that is associated with a nonsingular uncurling metric implies the full group, since the member implies the operation $s\mapsto s^{-1}$ via (\ref{11/7/20.1}), from which the full characterization group can be computed.  Note that this group is {\it not} $\mathfrak{u}\oplus U$.

The inversion characterization group of an algebra refers to characterizations of the algebra's multiplicative inversion operation $s\mapsto s^{-1}$, and does not necessarily fully characterize the algebra itself.  
Nevertheless,  the inversion characterization group embodies both a class of  novel algebra isomorphism invariants and (as shown in Section 3) a generalization of the usual algebra norm.  

Corollary \ref{8/18/23.2} has already identified invariants associated with $\mathfrak{u}$ (whose members supply the first component of the inversion characterization group), but another set of useful invariants becomes apparent on examination of $U$ (whose members supply the second component of the inversion characterization group).  Based on the mapping from $\mathfrak{u}$ to $ U$ defined by (\ref{11/7/20.2}),  we can write \begin{equation}  U = \exp\left(\int_{\bf 1}^s [\mathfrak{u}t^{-1}]\cdot dt\right). \label{6/7/23.1}\end{equation}  In the above, the notation $\mathfrak{u}t^{-1}$ means $\{Lt^{-1}: L\in \mathfrak{u}\}$.  Accordingly, the right-hand-side of (\ref{6/7/23.1}) is understood to be a collection of exponentiated integrals.   Note that in (\ref{6/7/23.1}) we do not need the factor $\frac{1}{\|{\bf 1}\|^2}$, since $\frac{1}{\|{\bf 1}\|^2}\mathfrak{u}=\mathfrak{u}$.  

\enlargethispage{2cm}



The invariants stemming from the $U$ contribution to the inversion characterization group proceed from the following. 



\begin{lemma}\label{6/6/23.2} Let $K:A_1\rightarrow A_2$ be an algebra isomorphism, where for $s_1\in A_1$ we define $s_2\equiv Ks_1$.  Then \[ \int_{{\bf 1}_{A_1}}^{s_1} \left[\mathfrak{u}_{A_1}t_1^{-1}\right]\cdot dt_1 
= \int_{{\bf 1}_{A_2}}^{s_2} \left[\mathfrak{u}_{A_2}t_2^{-1}\right]\cdot dt_2 . \] \end{lemma}

\begin{proof} 
Equation (\ref{5/23/23.1}) implies,
\begin{eqnarray*} \int_{{\bf 1}_{A_1}}^{s_1} \left[\mathfrak{u}_{A_1}t_1^{-1}\right]\cdot dt_1 &=& \int_{{\bf 1}_{A_1}}^{s_1} \left[ K^T\mathfrak{u}_{A_2}Kt_1^{-1}\right]\cdot dt_1 \\ &=& 
\int_{{\bf 1}_{A_2}}^{s_2} \left[ K^T\mathfrak{u}_{A_2}K (K^{-1}t_2)^{-1}\right]\cdot dK^{-1}t_2 \\ &=& \int_{{\bf 1}_{A_2}}^{s_2} [K^T\mathfrak{u}_{A_2}t_2^{-1}]\cdot \left(K^{-1}dt_2\right) = \int_{{\bf 1}_{A_2}}^{s_2} \left[\mathfrak{u}_{A_2}t_2^{-1}\right]\cdot dt_2, 
\end{eqnarray*}
where we have used $(K^{-1}t_2)^{-1} = K^{-1}t_2^{-1}$ (i.e., $\mathbf{1}_{A_1} = K^{-1}\mathbf{1}_{A_2} = K^{-1}(t_2t_2^{-1}) = (K^{-1}t_2)(K^{-1}t_2^{-1})$), the commutativity of a linear transformation and a differential, and exploitation of an adjoint in the context of an inner product.
\end{proof}

\begin{remark} {\rm As with Theorem \ref{6/6/23.1}, the above lemma easily generalizes to the case of uncurling metric groups of arbitrary oder.}\end{remark}

Applying Lemma \ref{6/6/23.2} to (\ref{6/7/23.1}) it is clear that for isomorphic algebras $A_1,A_2$, the groups $ U_{A_1},  U_{A_2}$ ``look the same" after accounting for the effect of the isomorphism on the vector space of elements of the algebras.

More specifically, with respect to some basis, we have noted that $\mathfrak{u}$ can be written as a matrix whose entries are each a linear combination of the components of an $m$-vector of real parameters $(\alpha_1, \dots, \alpha_m)$.   Thus, $\mathfrak{u}$ can be equivalently thought of as a set of real matrices, each of whose members is given by a choice of real values for the $m$ parameters in the above parametrized matrix version of $\mathfrak{u}$, with each matrix in the set being an uncurling metric.   Note that $\mathfrak{u}$ is an $m$-dimensional space.    Now, for each $i\le m$ consider the real matrix  $\mathfrak{u}_i$ obtained from $\mathfrak{u}$ by setting $\alpha_i =1$ with all other parameters set to zero.  These are a set of generators for $\mathfrak{u}$, given that the latter is also a real vector space. Thus, $\mathfrak{u}t^{-1} = \sum_{i=1}^m \alpha_i \mathfrak{u}_i t^{-1}$.   Each component of $\mathfrak{u}_i t^{-1}$ is a rational function, since the inverse of the left regular representation matrix associated with $t$ is the reciprocal of the determinant of that matrix multiplied by its adjugate matrix.  It follows that  with respect to a piecewise smooth path from $\mathbf{1}$ to $s$ such that $t$ only varies with respect to a single one of its components on any of the constituent smooth pieces, $\alpha_i \int_{\mathbf{1}}^s [\mathfrak{u}_i t^{-1}]\cdot dt$ can be treated by the method of partial fractions - so that the latter expression is seen to be the product of $\alpha_i$ with zero or one ``bare" rational function of the components of $s$, plus zero or one logarithm of a rational function of the components of $s$, plus zero or some positive integer number of arctangents of rational functions of the components of $s$.  We can notate the integral on the right-hand-side of (\ref{6/7/23.1}) as $\log  U$ (an $m$-dimensional collection of functions of $s$), which inherits a format as a single parametrized function from the parametrized format of $\mathfrak{u}$.   With respect to the chosen basis, we can write \begin{equation}\label{9/26/23.1} \log  U =\sum_{i=1}^m \alpha_i \int_{\mathbf{1}}^s [\mathfrak{u}_i t^{-1}]\cdot dt.\end{equation}   Invoking Lemma \ref{6/6/23.2}, we have   
\enlargethispage{1cm}
\begin{theorem} \label{6/6/23.3} Suppose $A_1$ and $A_2$ are isomorphic, and $\log U_{A_1}$ and $\log U_{A_2}$ are expressed according to (\ref{9/26/23.1}) following evaluation of the integrals.  Then these two parametrized functions have
\begin{itemize} \item  the same number of parameters multiplying logarithm terms,  \item  the same number of parameters multiplying arctangent terms, 
 \item   the same number of parameters multiplying bare rational function terms. \end{itemize}
\end{theorem}


So, for example, an alternative system of algebra classification  could be such that two algebras are in the same class if they have the same vector space dimension, they share the same values for the first three of the dimensional invariants of $\mathfrak{u}$ in the itemized list of Corollary \ref{8/18/23.2} (whose computations are straightforward), and they have the same respective numbers of parameters multiplying logarithm terms, arctangent terms, and bare rational function terms, as identified in Theorem \ref{6/6/23.3}.   The latter three numbers forwarded by Theorem \ref{6/6/23.3} constitute an invariant ``triple index" arising from the mapping from $\mathfrak{u}$ to $ U$ defined by (\ref{11/7/20.2}).  
Of course, one can imagine augmenting the class definition by introducing additional computationally straightforward algebra invariants as may be identified.  

\subsection{Illustrations using the two-dimensional algebras}
Up to an isomorphism, there are only three two-dimensional algebras.

For the algebra of complex numbers, $\mathbb{C}$, we can write its elements $s$ as $(x,y)$, with the product given by $(x,y)(z,w) = (xz-yw, yz+xw)$, with $\mathbf{1}=(1,0)$, so that $\|\mathbf{1}\|^2=1$.  Also, $(x,y)^{-1} = (\frac{x}{x^2+y^2},\frac{-y}{x^2+y^2})$.  According to the program indicated immediately prior to Corollary \ref{8/18/23.2}, we start out by writing  $\mathfrak{u}_\mathbb{C} = \left[\begin{array}{ccc}
 \alpha & \beta   \\ 
 \beta & \gamma    
\end{array}\right]$.  We then obtain  $\mathfrak{u}_\mathbb{C}s^{-1} = \left[\begin{array}{ccc}
 \alpha & \beta   \\ 
 \beta & \gamma    
\end{array}\right]\left(\begin{array}{ccc}
x   \\ 
y   
\end{array}\right)^{-1} = \left(\begin{array}{ccc}
\frac{\alpha x - \beta y}{x^2+y^2}  \\ 
\frac{\beta x - \gamma y}{x^2+y^2} 
\end{array}\right)$.  As regards the right-hand-side of this equation, satisfaction of (\ref{4/7/23.1}) requires that the partial derivative with respect to $x$ of the second component must equal the partial derivative with respect to $y$ of the first component.  Consequently, $\mathfrak{u}_{\mathbb{C}} = \left[\begin{array}{ccc}
 \alpha & \beta   \\ 
 \beta & -\alpha.    
\end{array}\right]$.  A simple exercise in integration leads to $\log  U_{\mathbb{C}} = \frac{\alpha}{2}\log(x^2+y^2) + \beta{\rm arctan}(\frac{y}{x})$, an expression relevant to Theorem \ref{6/6/23.3}.

The split-complex numbers $\slashed{\mathbb{C}}$ is an algebra for which the elements are  $(x,y)$ with product $(x,y)(z,w) = (xz+yw,xw+yz)$, $ \mathbf{1} = (1,0)$, and $\|\mathbf{1}\|^2 = 1$.  Also, $(x,y)^{-1} = (\frac{x}{x^2-y^2},\frac{-y}{x^2-y^2})$.  We then find $\mathfrak{u}_{\slashed{\mathbb{C}}} = \left[\begin{array}{ccc}
 \alpha & \beta   \\ 
 \beta & \alpha.    
\end{array}\right]$, and $\log  U_{\slashed{\mathbb{C}}} = \frac{\alpha+\beta}{2}\log(x+y) + \frac{\alpha-\beta}{2}\log(x-y) = \frac{\alpha}{2}\log[(x+
y)(x-y)] +  \frac{\beta}{2}\log \frac{x+y}{x-y}$.

The final real two-dimensional unital algebra (up to an isomorphism) is the algebra of dual numbers $\mathbb{D}$.  The elements are $(x,y)$ with product given by $(x,y)(z,w) = (xz,xw+yz)$, with $\mathbf{1}=(1,0)$ and $\|\mathbf{1}\|^2 = 1$. Also, $(x,y)^{-1} = (\frac{1}{x},-\frac{y}{x^2})$.  We then find $\mathfrak{u}_\mathbb{D} = \left[\begin{array}{ccc}
 \alpha & \beta   \\ 
 \beta & 0    
\end{array}\right]$ and $\log  U_{\mathbb{D}} = \alpha\log x + \beta\left(\frac{y}{x}\right)$.

Now we consider $\mathbb{R}\oplus\mathbb{R}$, which is isomorphic to $\slashed{\mathbb{C}}$.  We have elements as $(x,y)$, product $(x,y)(z,w) = (xz,yw)$, $\mathbf{1}=(1,1)$, and $\|\mathbf{1}\|^2 = 2$.  Also $(x,y)^{-1} = (\frac{1}{x},\frac{1}{y})$.  We then find $\mathfrak{u}_{\mathbb{R}\oplus\mathbb{R}} = \left[\begin{array}{ccc}
 \alpha & 0   \\ 
 0 & \gamma    
\end{array}\right]$ and $\log  U_{\mathbb{R}\oplus\mathbb{R}} = \frac{\alpha}{2}\log x + \frac{\gamma}{2}\log y$.

We can observe that the parametrized function formats associated with the respective $\log U_{\mathbb{C}}$, $\log U_{\slashed{\mathbb{C}}}$, and $\log U_{\mathbb{D}}$, are incompatible with any pair of the associated algebras being isomorphic according to Theorem \ref{6/6/23.3}.  However, the function formats appearing in the respective $\log U_{\slashed{\mathbb{C}}}$ and $\log U_{\mathbb{R}\oplus\mathbb{R}}$  for the isomorphic $\slashed{\mathbb{C}}$ and $\mathbb{R}\oplus\mathbb{R}$ are entirely compatible.  This is an illustration of use of isomorphism invariants arising from $U$.

  As regards the isomorphism invariants stemming from $\mathfrak{u}$, the dimension of $\mathfrak{u}$ is $2$ for each of these algebras.  But the dimension of the subspace of $\mathfrak{u}$ to which the determinant is sensitive is $2$, $2$, and $1$, for $\mathbb{C}$, $\slashed{\mathbb{C}}$, and $\mathbb{D}$, respectively.  The real dimension of the real algebraic set ${\rm det}(M_{\mathfrak{u}}) = 0$ is $0$, $1$, and $1$, for $\mathbb{C}$, $\slashed{\mathbb{C}}$, and $\mathbb{D}$, respectively.  This indicates the utility of isomorphism invariants arising from $\mathfrak{u}$. 

\subsection{Special unital norms}

On a convex neighborhood of $\mathbf{1}$ devoid of non-units but containing $\left(\frac{1}{\alpha}\right)\mathbf{1}$ for $\alpha>0$, it is easy to show from (\ref{11/7/20.2}) that for unital norms in general, \[  \ell(\alpha s) = \alpha^{\frac{\mathbf{1}^T L \mathbf{1}}{\|\mathbf{1}\|^2}} \ell(s),\] using an argument similar to that employed in the proof of Theorem \ref{9/24/23.1}. Thus, $ \ell(s)$ is a positive homogeneous function of degree given by the exponent on the right-hand-side of the above equation.  
   {\it But things become somewhat nicer if elements that are inverses with respect to the algebra's product also behave inversely with respect to the scalar product defined by $L$}.  
    This situation results if uncurling metric $L$ satisfies the additional constraint (\ref{7/7/21.1}).
 
\begin{definition}\label{9/20/23.1} An uncurling metric $L$ is a {\it normalized uncurling metric} if for all units it satisfies, \begin{equation} \label{7/7/21.1} s^TLs^{-1} =\|{\bf 1}\|^2.\end{equation}  A {\it special unital norm} is a unital norm $ \ell(s)$ which is derived from a normalized uncurling metric $L$ according to (\ref{11/7/20.2}). 
\end{definition}

\begin{theorem}\label{9/24/23.1} If $ \ell(s)$ is a special unital norm then it is a degree-1 positive homogeneous function, and
\begin{equation}\label{5/23/23.3}  \ell(s^{-1}) = \big(\ell(s)\big)^{-1} = \frac{1}{\ell(s)}, \end{equation} for $s,s^{-1}$ in some convex neighborhood of $\mathbf{1}$ devoid of non-units. \end{theorem}

\begin{proof} We have already observed (following Definition \ref{9/21/23.2}) that a convex neighborhood of $\mathbf{1}$ devoid of non-units always exists.  Let $\mathcal{N}$  be such a neighborhood.

For $s,\alpha s\in\mathcal{N}$,  and the line segment from $\mathbf{1}$ to $\frac{\mathbf{1}}{\alpha}\mathbf{1}$ also in  $\mathcal{N}$, with $\alpha>0$, we have from (\ref{11/7/20.2}) that
\begin{eqnarray}\nonumber  \log\ell(\alpha s) &=&  \frac{1}{\|{\bf 1}\|^2}\int_{\bf 1}^{\alpha s} [Lt^{-1}]\cdot dt = \frac{1}{\|{\bf 1}\|^2}\int_{\frac{\mathbf{1}}{\alpha}\mathbf{1}}^{s} [L(\alpha w)^{-1}]\cdot d(\alpha w)  \\ \nonumber &=&  \frac{1}{\|{\bf 1}\|^2}\int_{\mathbf{1}}^{s} [Lw^{-1}]\cdot dw + \frac{1}{\|{\bf 1}\|^2}\int_{\frac{\mathbf{1}}{\alpha} \mathbf{1}}^{\mathbf{1}} [Lw^{-1}]\cdot dw  \\ &=& \log \ell(s) + \log\alpha, \label{4/2/23.1}
\end{eqnarray} where the integral from $\frac{\mathbf{1}}{\alpha} \mathbf{1}$ to $\mathbf{1}$ can be easily evaluated along the line segment with those endpoints, and we have used (\ref{7/7/21.1}) in evaluation of that integral to obtain the term $\log \alpha$ on the right-hand-side of the final equation.  Thus, we have $\ell(\alpha s) = \alpha\ell(s)$, degree-1 positive homogeneity.

Now consider a smooth path $\mathbf{P}$ from $\mathbf{1}$ to $s$ which is inside $\mathcal{N}$  such that $\mathbf{P}^{-1}$ is also inside $\mathcal{N}$,  where $\mathbf{P}^{-1}$ is the point set consisting of the multiplicative inverses of the members of $\mathbf{P}$.  The existence of neighborhood $\mathcal{N}$ with such a property is easily demonstrated using the left regular representation of the algebra.
\mbox{We then have,} \begin{eqnarray}  \log \ell(s^{-1})   &=&   \frac{1}{\|{\bf 1}\|^2} \int_{\bf 1}^{s^{-1}} [Lt^{-1}] \cdot dt \nonumber \\ &=&  \frac{1}{\|{\bf 1}\|^2} \int_{\bf 1}^s [Ly]\cdot d\left(y^{-1}\right) \nonumber \\ &=&  \frac{1}{\|{\bf 1}\|^2}\left(\left(y^{-1}\cdot [Ly]\right)\bigg|_{\bf 1}^s - \int_{\bf 1}^s y^{-1}\cdot d\left(Ly\right)\right)   \nonumber \\ &=& -\frac{1}{\|{\bf 1}\|^2} \int_{\bf 1}^s y^{-1}\cdot (Ldy) = -\frac{1}{\|{\bf 1}\|^2} \int_{\bf 1}^s [Ly^{-1}] \cdot dy  \nonumber \\ &=& -\log \ell(s),\label{4/2/23.5} \end{eqnarray} where we have used a change of variable (contemplating the left regular representation, it is appreciated that $y=t^{-1}$ is a diffeomorphism on $\mathcal{N}$), commutativity of a linear transformation and a differential, the property that $L$ is a real symmetric matrix, and  (\ref{7/7/21.1}).  Equation (\ref{5/23/23.3}) then follows.
\end{proof}
As demonstrated in Section 3, every algebra is associated with at least one special unital norm, which ultimately follows from the fact that the transpose operation is a special unital norm for $M_n(\mathbb{R})$.

It is also seen in Section 3 that at least some members of an algebra's group of unital norms are multiplicative, i.e., $\ell(s_1s_2) = \ell(s_1)\ell(s_2)$ (in fact, all of the unital norms associated with the two-dimensional algebras are multiplicative).   But in general, unital norms are not multiplicative.  However, the unital norms originate from the nonlinear unary multiplicative inverse operation $s\mapsto s^{-1}$ implied by the algebra's bilinear binary product operation - so the particularly desirable fealty would be to that unary operation.  It is thus significant that the special unital norms pay homage to it via (\ref{5/23/23.3}) (also, see comments in the second and third paragraphs of Section 4.2). 

\subsection{The functor}

We have been referring to algebra isomorphisms, but this approach can also more generally address algebra epimorphisms, via a particular functor introduced as follows.  

Let $\mathcal{A}$ be the category each of whose objects is an algebra (i.e., a real finite-dimensional unital associative algebra), and whose morphisms are algebra epimorphisms (we will use the same symbol, e.g., $K$, for both the epimorphism and morphism).  We define category $\mathcal{P}$ as follows. An object in this category is a symmetric $(n\times n)$-matrix where $n$ is any particular positive integer, and where each matrix entry is a linear combination of the components of an $m$-vector of real parameters, $m\ge 1$ (i.e., the $m$ components of the later vector are each a different real parameter).  Thus, for any object of $\mathcal{P}$, the implied set of real matrices resulting from all realizations of the associated $m$-vector of parameters is an $m$-dimensional space (a realization of the $m$-vector of parameters is a real $m$-vector resulting from a choice of real values for the $m$ parameters).  An object of $\mathcal{P}$ is considered to be the same as the latter space of real matrices.  Consider two objects ${\sf P}_1,{\sf P}_2$ of $\mathcal{P}$, where ${\sf P}_1$ is a $(j\times j)$-matrix parametrized by the components of a $k$-parameter vector, and ${\sf P}_2$ is a $(q\times q)$-matrix parametrized by the components of a $r$-parameter vector.  There is a morphism $ \begin{CD}  {\sf P}_1 @>p_1>> {\sf P}_2\end{CD}$ if and only if  there exists a real $(q\times j)$-matrix $V$ such that $VV^T$ is nonsingular and $V^T{\sf P}_2V\subset {\sf P}_1$, where the latter subset expression is 
interpreted to mean that the set of real matrices implied by $V^T{\sf P}_2V$ through all realizations of the $r$-parameter vector associated with ${\sf P}_2$, is a subset of  the set of real matrices implied by  ${\sf P}_1$ through all realizations of the associated $k$-parameter vector. It is evident that there is at most one morphism pointing from one category $\mathcal{P}$ object to another.  It is also easily seen that the morphisms are associative.  

We define a mapping $\mathbf{F}$ from category $\mathcal{A}$ to category $\mathcal{P}$ as follows.  $\mathbf{F}[
A] \equiv \mathfrak{u}_A$, and given a morphism $K$ from $A_1$ to $A_2$, we define $\mathbf{F}(K)$ to be the morphism from $\mathbf{F}[A_1]$ to $\mathbf{F}[A_2]$, if the morphism exists.   

\begin{theorem}\label{6/13/22.1}  ${\bf F}: \mathcal{A}\rightarrow \mathcal{P}$, is a covariant functor. 
 \end{theorem}
 \begin{proof}
We will first show that the category $\mathcal{A}$ morphism $ \begin{CD}  A_1 @>K>> A_2\end{CD}$ implies the existence of category $\mathcal{P}$ morphism $ \begin{CD}  {\bf F}[A_1] @>{\bf F}(K)>> {\bf F}[A_2]\end{CD}$.

Morphism $K$ results from an algebra epimorphism $K:A_1\rightarrow A_2$.   Taking $s_1$ to be notation for a variable in $A_1$, we have ${\bf 1}_{A_2} = K(s_1s_1^{-1}) = K(s_1)K(s_1^{-1})$, so that $(Ks_1)^{-1} = K(s_1^{-1})$.  By definition, given the uncurling metric group $\mathbf{F}[A_2]$ associated with $A_2$ and an arbitrary real matrix  $\hat{\mathbf{F}}[A_2]$ in the set of real matrices implied by $\mathbf{F}[A_2]$, we have $d\left((s_2^{-1})^T\hat{\mathbf{F}}[A_2] (ds_2)\right) = 0$ in some neighborhood of ${\bf 1}_{A_2}$, where $s_2$ is notation for a variable in $A_2$.   With respect to appropriate neighborhoods of $\mathbf{1}_{A_1}$ and $\mathbf{1}_{A_2}$, it follows that \begin{eqnarray}\nonumber 0 = d\left((s_2^{-1})^T\hat{\mathbf{F}}[A_2] ds_2\right) &=&  d\left(((Ks_1)^{-1})^T\hat{\mathbf{F}}[A_2] d(Ks_1)\right) \\  &=& d\left((s_1^{-1})^T \big[K^T \hat{\mathbf{F}}[A_2]K\big]ds_1\right), \label{7/4/22.5}  \end{eqnarray}   using the equation at the end of the second sentence of this paragraph, and the fact that a differential commutes with a linear transformation.  Equation (\ref{7/4/22.5}) is simply an instance of (\ref{4/7/23.1}), so it follows that symmetric $K^T \hat{\mathbf{F}}[A_2]K$ is an uncurling metric associated with $A_1$.  Consequently, $K^T\mathbf{F}[A_2]K\subset \mathbf{F}[A_1]$.  Furthermore, $KK^T$ is nonsingular since $K$ is an epimorphism.  Hence, criteria for the existence of morphism $\mathbf{F}(K)$ are fulfilled (with $V=K)$.

Commutativity of the requisite diagram
  \[  \begin{CD}  
A_1 @>K_1>> A_2 @>K_2>> A_3\\
 @VVV   @VVV  @VVV\\
 {\bf F}[A_1]   @>{\bf F}(K_1)>>  {\bf F}[A_2]  @>{\bf F}(K_2)>>  {\bf F}[A_3] \\ 
 \end{CD} \]
is then evident from existence and uniqueness of the relevant category $\mathcal{P}$ morphisms, and the mechanism of their composition.  Thus, ${\bf F}$ is a covariant functor. 
\end{proof}

We then have the following useful results.  The first is immediate.

\begin{corollary}\label{6/26/22.2} An epimorphism from one algebra to another algebra can exist only if there is a category $\mathcal{P}$ morphism associated with their respective uncurling metric groups. \end{corollary}

\begin{corollary} \label{8/17/22.5} An epimorphism from a first algebra to a second algebra can exist only if the dimension of the uncurling metric group of the first algebra is greater than or equal to the dimension of the uncurling metric group of the second algebra. \end{corollary}  

\begin{proof} Given algebra epimorphism $K:A_1\rightarrow A_2$, suppose $A_1$ is a $j$-dimensional algebra with an uncurling metric group utilizing a $k$-parameter vector (and thereby the uncurling metric group is a $k$-dimensional space), and $A_2$ is a $q$-dimensional algebra with an uncurling metric group utilizing a $r$-parameter vector (and thereby the uncurling metric group is a $r$-dimensional space).  Evidently, $K$ is a $(q\times j)$-matrix with $q\le j$ and $KK^T$ is nonsingular (since $K$ is an epimorphism). The proof of Theorem \ref{6/13/22.1} has shown that $K^T \mathbf{F}[A_2] K \subset \mathbf{F}[A_1]$.  Thus, $K^T \mathbf{F}[A_2] K$ cannot have more parameters than $\mathbf{F}[A_1]$  (i.e., the dimension of $K^T \mathbf{F}[A_2] K$ cannot exceed the dimension of $\mathbf{F}[A_1]$).  Furthermore, $K^T \mathbf{F}[A_2] K$ has at least as many parameters as $KK^T \mathbf{F}[A_2] KK^T$, but it has less than or equal to the number of parameters of $\mathbf{F}[A_2]$.  However, $KK^T \mathbf{F}[A_2] KK^T$ has the same number of parameters as $\mathbf{F}[A_2]$ since $KK^T$ is nonsingular.  Thus, $K^T \mathbf{F}[A_2] K$ has the same number of parameters as $\mathbf{F}[A_2]$.  Hence, the dimension of $\mathbf{F}[A_2]$ cannot exceed the dimension of $\mathbf{F}[A_1]$.
\end{proof}

\begin{corollary}\label{8/17/22.1} If the uncurling metric group of an algebra is one-dimensional and contains a nonsingular member, then the algebra is simple. \end{corollary}

\begin{proof} Suppose algebra $A_1$ is $j$-dimensional and has a non-trivial ideal, and its uncurling metric group is one-dimensional and has a nonsingular member.  It follows that $\mathbf{F}[A_1]$ is represented by  $\alpha L_1$, $\alpha\in \mathbb{R}$, where $L_1$ is a real nonsingular $(j\times j)$-matrix.  Since $A_1$ has a non-trivial ideal, there exists an algebra epimorphism $K:A_1\rightarrow A_2$ where $A_2$ has dimension $q<j$.  Corollary \ref{8/17/22.5} indicates that the uncurling metric group of $A_2$ cannot have dimension greater than one, and thus $\mathbf{F}[A_2]$ is represented by $\alpha L_2$ where $L_2$ is a real $(q\times q)$-matrix with $q<j$. From the proof of Theorem \ref{6/13/22.1}, we must then have that the nonzero matrix $ K^T L_2 K$ is a member of the uncurling metric group $\mathbf{F}[A_1]$.  But the determinant of the nonzero matrix $ K^T L_2 K$ is zero since $K$ has a nontrivial kernel, while the only member of the uncurling metric group $\mathbf{F}[A_1]$ having zero determinant is the zero matrix (corresponding to $\alpha=0$), since $L_1$ is necessarily nonsingular.   It follows that $K$ cannot exist, and so $A_1$ is simple. \end{proof}


In view of Theorem \ref{4/8/22.1}, an easy consequence of Corollary \ref{8/17/22.1} is an alternative demonstration that the algebra of real $(n\times n)$-matrices $ M_n(\mathbb{R})$ is simple.  

Note that the converse of Corollary \ref{8/17/22.1} is not true.  For example, $\mathbb{C}$ is simple but has a two-dimensional uncurling metric group (see Section 2.3).


\section{Unital norms generalize the usual algebra norm} 
The ``usual" algebra norm is given by the determinant of an element's image under the left regular representation (e.g., \cite{bourbaki:1989}).  The group $U$ of unital norms subsumes this definition in the sense of the theorems and corollaries of this section.  This also means that the uncurling metric group of an algebra is at least one-dimensional.

 We begin with the matrix algebra $ M_n(\mathbb{R})$. For an invertible element $s$, with det$(s)$ being its determinant, we have the well-known
  \begin{equation}\label{11/5/21.1} \nabla \mbox{det}(s) = \big[{\rm Adj}(s)\big]^T = [s^{-1}]^T\mbox{det}(s),\end{equation} where Adj$(s)$ is the adjugate matrix.  The first of these equalities is essentially Jacobi's formula, and the second follows from the classical matrix inverse expression resulting from Laplace expansion.  The gradient notation on the left-hand-side is employed with the understanding that an element of the algebra (including the full expression on the right-hand-side), as a member of the vector space of elements, is identified as the $n^2$-tuple obtained from sequentially appending the columns of its matrix format (and vice-versa), with a compatible understanding regarding the transpose operation as a particular involutive $(n^2\times n^2)$-matrix, which we will denote as $\mathcal{T}$. 
  
From (\ref{11/5/21.1}), $\mathcal{T}s^{-1} = \frac{\nabla{\rm det}(s)}{{\rm det}(s)} = \nabla[\log{\rm det}(s)]$, so $\mathcal{T}$ is an uncurling metric.  Since ${\rm det}(s)$ is a degree-$n$ positive homogeneous function, the Euler Homogeneous Function Theorem implies $[\nabla{\rm det}(s)]\cdot s = n \,{\rm det}(s)$.  We also have $\|\mathbf{1}\|^2 = n$.  Consequently, (\ref{7/7/21.1}) is satisfied for $L=\mathcal{T}$, so $\mathcal{T}$ is a normalized uncurling metric.  Applying the first equation of this paragraph to (\ref{11/7/20.2}), it is then seen that $(\text{det}(s))^{\frac{1}{n}}$ is a special unital norm.
  Moreover,
 
   
 \begin{theorem} \label{4/8/22.1} $\mathfrak{u}_{M_n(\mathbb{R})}$ is one-dimensional and there is only a single normalized uncurling metric, this being the  transpose operation on $(n\times n)$-matrices expressed as an $(n^2\times n^2)$-matrix. The associated special unital norm is $({\rm det}(s))^{\frac{1}{n}}$. \end{theorem}
 \begin{proof} The vector space of members of $ M_n(\mathbb{R})$ consists of  vectors with $n^2$ components.   Using the sequentially appended matrix columns convention as introduced in the second paragraph of this section, we can equivalently represent a member $s\in M_n(\mathbb{R})$ in either matrix or column vector format.   An uncurling metric for this algebra will be a matrix $L\in M_{n^2}(\mathbb{R})$, and must satisfy (\ref{4/7/23.1}), i.e., $d\left([Ls^{-1}]\cdot ds\right)=0$. This constitutes a condition on every pair of distinct components of $Ls^{-1}$ (the latter having $n^2$ components).  Thus,  the partial derivative of any component of $Ls^{-1}$, e.g., the component corresponding to the $(i,j)$ entry of matrix $s$, taken with respect to, e.g., the $(k,m)$ component of matrix $s$ ($(k,m)\ne (i,j)$),  minus  the partial derivative of the $(k,m)$ component of $Ls^{-1}$  taken with respect to the $(i,j)$ component of matrix $s$, must be zero.
 
 To examine this constraint, we recall that Matrix Calculus supplies the derivative of the inverse of a matrix, $\frac{\partial (s^{-1})}{\partial u}$ (with $u$ a matrix), as follows.  Differentiating both sides of the expression
$ ss^{-1} = {\bf 1}$, we obtain $\frac{\partial (s^{-1}s)}{\partial u} = \frac{\partial (s^{-1})}{\partial u}s + s^{-1}\frac{\partial s}{\partial u}=0$ (where $0$ is the matrix of all zero entries).  Hence, 
 \[ \frac{\partial (s^{-1})}{\partial u} = -s^{-1}\frac{\partial s}{\partial u}s^{-1},\] which is an $(n\times n)$-matrix.  Let $\{ e_{pq}\}$ be the canonical basis for $ M_n(\mathbb{R})$, so that a basis member $e_{rw}$ is the matrix whose $(r,w)$ entry is $1$, but with all other entries zero.  It follows that $-s^{-1}\frac{\partial s}{\partial e_{pq}}s^{-1} = -s^{-1} e_{pq}s^{-1}$  is a matrix whose entries are the $n^2$ components of the derivative of $s^{-1}$ in the $e_{pq}$ direction (i.e., the full derivative of $s^{-1}$ is an array of all such matrices in each of the basis members of $\{ e_{pq}\}$).  
 Thus, the $e_{km}$-direction derivative matrix of $s^{-1}$ is, \begin{equation}\label{8/18/22.5} -s^{-1} e_{km}s^{-1} = -(s^{-1})_{:k} (s^{-1})_{m:},\end{equation} where $(s^{-1})_{:m}$ indicates the $m$-th column of $s^{-1}$, and $(s^{-1})_{k:}$ indicates the $k$-th row of $s^{-1}$.    Similarly, the $e_{ij}$-direction derivative matrix of $s^{-1}$ is \begin{equation}\label{8/18/22.6} -s^{-1} e_{ij}s^{-1} = -(s^{-1})_{:i} (s^{-1})_{j:}.\end{equation}  It is thus clear that the sets of entries of the $e_{km}$-direction derivative matrix and the $e_{ij}$-direction derivative  matrix share the member  $-(s^{-1})_{mi} (s^{-1})_{jk}$, and they can share no other member because $\{(s^{-1})_{pq}\} = \{\frac{C_{qp}}{\det(s)}\}$, where the right-hand-side of this equation is clearly a set of linearly independent functions since $\{C_{qp}\}$ is the matrix of cofactors of $s$. The latter observation also indicates that the union of the entries of the $e_{km}$-direction derivative matrix with the entries of the $e_{ij}$-direction derivative matrix is a set of linearly independent functions.

Therefore, the partial derivative in the $e_{km}$-direction of the $e_{ij}$-component of $Ls^{-1}$ is a linear combination of the entries of the matrix on the right-hand-side of (\ref{8/18/22.5}), and  the partial derivative in the $e_{ij}$-direction of the $e_{km}$-component of $Ls^{-1}$ is a linear combination of the entries on the right-hand-side of (\ref{8/18/22.6}).
So satisfaction of $d\left([Ls^{-1}]\cdot ds\right) = 0$, based on the observation made at the end of the prior paragraph, requires that all entries in any two rows of $L$ must be zero, except for the one entry in each row corresponding to the shared member of the two matrix set of entry members of the two relevant partial derivative matrices (the right-hand-sides of (\ref{8/18/22.5}) and (\ref{8/18/22.6})) noted in the prior paragraph - and these two corresponding nonzero entries of the two cited rows of $L$ must have the same value.  This prescription must be true for any two rows of $L$, so all nonzero entries of $L$ must have the same value, and there can be at most one nonzero value in each row of $L$.  Thus, the uncurling metric group is at most one-dimensional. Assuming the group is one-dimensional, the uncurling metric group can be represented as the product of a constant matrix with a scalar $\alpha$ representing the uniform value of the nonzero entries of $L$.  
For the special unital norm, the value of $\alpha$ is determined by the requirement that (\ref{7/7/21.1}) be satisfied. 

Hence, there is at most one possible normalized uncurling metric.  But according to observations made in the third paragraph of this section, $ M_n(\mathbb{R})$ does indeed have a normalized uncurling metric - this being the $(n^2\times n^2)$-matrix corresponding to the matrix transpose operation on $(n\times n)$-matrices.  As was shown immediately prior to the theorem statement, $(\mbox{det}(s))^{\frac{1}{n}}$ is the associated special unital norm.  
 \end{proof}
 
 The following theorem and corollary indicate the manner in which unital norms generalize the usual algebra norm.
 
 \begin{theorem}\label{5/13/22.1} Let $\bar{A}$ be an $n$-dimensional algebra and let $A$ be its left regular representation.  If $A$ is such that the transpose operation on $ M_n(\mathbb{R})$ maps $A$ to itself, then the special unital norm family of $\bar{A}$ has a member $\ell$ such that  for $\bar{s}\in\bar{A}$ and $s$ its image under the left regular representation, $\ell(\bar{s}) = \left({\rm det}(s)\right)^{\frac{1}{n}}$. 
\end{theorem}

\begin{proof} 
Algebra $\bar{A}$ is isomorphic to its left regular representation $A$, which is a subalgebra  of $ M_n(\mathbb{R})$.   $\mathbf{1}_{M_n(\mathbb{R})}$ is the $(n\times n)$-identity matrix, and observe that as an element of $ M_n(\mathbb{R})$, $\mathbf{1}_A$ is likewise the $(n\times n)$-identity matrix.  However, $ M_n(\mathbb{R})$ is $n^2$-dimensional with $\|\mathbf{1}\|^2 = n$, while $A$ is $n$-dimensional with $\|\mathbf{1}_A\|^2 = \|\mathbf{1}_{\bar{A}}\|^2$.  That is, as {\it vector spaces}, $\bar{A}$ is mapped to  $A$ by the identity operation, and the latter are subject to the Euclidean inner product as per the conventions stated in the Introduction.   

As noted shortly prior to the statement of Theorem \ref{4/8/22.1}, the transpose operation on $ M_n(\mathbb{R})$ is a linear transformation $\mathcal{\mathcal{T}}:\mathbb{R}^{n^2}\rightarrow \mathbb{R}^{n^2}$ and, also notating it as $\mathcal{\mathcal{T}}$ with respect to the standard basis, matrix $\mathcal{\mathcal{T}}$ represents a normalized uncurling metric associated with algebra  $ M_n(\mathbb{R})$.  Being an uncurling metric means that integration of $\mathcal{\mathcal{T}}s^{-1}$ is path-independent in some simply connected neighborhood of the multiplicative identity of $ M_n(\mathbb{R})$.  
 Hence, $\int_{\mathbf{1}_{M_n({\mathbb{R})}}}^s (\mathcal{\mathcal{T}}t^{-1})\cdot dt$ is path-independent in the latter neighborhood and, in particular, is path-independent for paths in that neighborhood restricted to be in $A$.  By assumption, $\mathcal{\mathcal{T}}$ maps $A$ to $A$. 
 Thus, the curl of $\mathcal{\mathcal{T}}s^{-1}$ vanishes when the domain of integration is restricted to $A$, i.e., 
  $d((\mathcal{\mathcal{T}}s^{-1})\cdot ds) = 0$ in some neighborhood of $\mathbf{1}_A$ in $A$, since this is a necessary condition for $\mathcal{T}$ to be an uncurling metric on $ M_n(\mathbb{R})$.
  
   Accordingly, we can use $\mathcal{\mathcal{T}}$ to define a scalar product on $A$ (not necessarily positive-definite or nondegenerate) as $\langle s_1,s_2\rangle_A \equiv  \frac{\|\mathbf{1}_{A}\|^2}{n}s_1^T \mathcal{\mathcal{T}} s_2$. 
    Since $A$ and $\bar{A}$ are isomorphic, the above scalar product on $A$ can be employed to produce a scalar product on $\bar{A}$ so that the two resulting quadratic spaces are isomorphic.  Specifically,  for $s_1,s_2\in A$, let $ \mathcal{I}:\bar{A}\rightarrow A$ be the left regular representation isomorphism, with $s_1 = \mathcal{I}\bar{s}_1$ and $s_2= \mathcal{I}\bar{s}_2 $.  Then we define the induced scalar product on $\bar{A}$ via, \begin{equation}\label{7/15/22.1} \bar{\mathcal{\mathcal{T}}} \equiv  \frac{\|\mathbf{1}_{\bar{A}}\|^2}{n}\mathcal{I}^T\mathcal{\mathcal{T}}\mathcal{I},\end{equation} i.e., $\langle \bar{s}_1,\bar{s}_2\rangle_{\bar{A}} \equiv \bar{s}_1^T \bar{\mathcal{\mathcal{T}}} \bar{s}_2$.  Applying (\ref{7/15/22.1}) and $\|\mathbf{1}_A\|^2 = \|\mathbf{1}_{\bar{A}}\|^2$, we have \begin{equation}\label{8/25/22.1} \langle \bar{s}_1,\bar{s}_2\rangle_{\bar{A}} = \bar{s}_1^T \bar{\mathcal{T}} \bar{s}_2 = \frac{\|\mathbf{1}_{\bar{A}}\|^2}{n}( s_1^T \mathcal{T} s_2) = \langle s_1,s_2\rangle_{A}.\end{equation}  The fraction $(\|\mathbf{1}_{\bar{A}}\|^2/n)$ is employed in (\ref{7/15/22.1}) in order that $\langle \mathbf{1}_{\bar{A}},\mathbf{1}_{\bar{A}}\rangle_{\bar{A}} = \|\mathbf{1}_{\bar{A}}\|^2$, given that we have $\mathbf{1}_{M_n(\mathbb{R})}^T \mathcal{T}\,\mathbf{1}_{M_n(\mathbb{R})} = n = \|\mathbf{1}_{M_n(\mathbb{R})}\|^2$. Due to this isomorphism of $A$ and $\bar{A}$ as quadratic spaces,  in small neighborhoods of the respective multiplicative identities $\mathbf{1}_A$ and $\mathbf{1}_{\bar{A}}$, the previously demonstrated $d((\mathcal{T}s^{-1})\cdot ds) = 0$ implies 
\mbox{$d((\bar{\mathcal{T}}(\bar{s})^{-1})\cdot d\bar{s}) = 0$.}. Thus, $\bar{\mathcal{T}}$ is an uncurling metric for $\bar{A}$.


In the context of $M_n(\mathbb{R})$, we know from Theorem \ref{4/8/22.1} that $L=\mathcal{T}$ satisfies (\ref{7/7/21.1}), where $\|\mathbf{1}\|^2$ in  (\ref{7/7/21.1}) is $\|\mathbf{1}_{M_n({\mathbb{R})}}\|^2 = n$.  Combining these facts with (\ref{8/25/22.1}), we have \begin{equation}\label{9/27/23.1} (\bar{s})^T \bar{\mathcal{T}} (\bar{s})^{-1} = \frac{\|\mathbf{1}_{\bar{A}}\|^2}{n}( s^T \mathcal{T} s^{-1}) = \|\mathbf{1}_{\bar{A}}\|^2 .\end{equation}
That is, $\mathcal{\bar{T}}$ is a normalized uncurling metric.  Hence, a special unital norm on $\bar{A}$ results from (\ref{11/7/20.2}) (with $\mathbf{1}=\mathbf{1}_{\bar{A}}$)  via the normalized uncurling metric $L = \bar{\mathcal{T}}$, as derived from the matrix transpose operation $\mathcal{T}$. 
 Again applying (\ref{8/25/22.1}) and (\ref{7/7/21.1}),  we then obtain 
\begin{eqnarray}\nonumber \log\ell(\bar{s}) = \frac{1}{\|\mathbf{1}_{\bar{A}}\|^2} \int_{\mathbf{1}_{\bar{A}}}^{\bar{s}} (\bar{\mathcal{T}}(\bar{s})^{-1})\cdot d\bar{s}
 &=& \frac{1}{\|\mathbf{1}_{\bar{A}}\|^2} \int_{\mathbf{1}_{A}}^{s} \frac{\|\mathbf{1}_{\bar{A}}\|^2}{n}(\mathcal{T}s^{-1})\cdot ds \\ &=& \frac{1}{n} \log\text{det}(s),\label{9/27/23.2}\end{eqnarray} where the final equation results from Theorem \ref{4/8/22.1}.
The theorem then follows. 
 \end{proof} 
 
 \newpage
 
 
 \begin{corollary}\label{11/17/22.1} Let $\bar{A}$ be an $n$-dimensional algebra, with $A$ its left regular representation, and $\mathcal{I}:\bar{A}\rightarrow A$ the isomorphism.  Let $\hat{A}$ be a subalgebra of $A$ having the property that the transpose operation $\mathcal{T}$ on $M_n(\mathbb{R})$ maps $\hat{A}$ to itself, and $\hat{A}$ is the largest subalgebra of $A$ with this property. Let $\underline{A}$ be the subalgebra of $\bar{A}$ given by application of $\mathcal{I}^{-1}$ to the members of $\hat{A}$, and let $\Omega$ be the operator yielding the orthogonal projection of $\bar{A}$ to $\underline{A}$. Then the special unital norm family associated with $\bar{A}$ has a member $\ell$ such that for $\bar{s}\in\bar{A}$ with $\hat{s}$ the component of $\mathcal{I}(\bar{s})$ that \mbox{is in $\hat{A}$,} \begin{equation}\label{2/11/24.1} \ell(\bar{s}) = \left({\rm det}(\hat{s})\right)^{\frac{1}{m}},\end{equation} where $m\equiv \mathbf{1}_{\bar{A}}^T(\Omega^T\mathcal{I}^T\mathcal{T}\,\mathcal{I}\,\Omega)\mathbf{1}_{\bar{A}}$. \end{corollary}
 
 \begin{proof} The proof proceeds along the same lines as the proof of Theorem \ref{5/13/22.1}.  The only significant difference is that, in reference to the integral in (\ref{11/7/20.2}), instead of projecting a path of integration in the vector space of $M_n(\mathbb{R})$ to a path in $\bar{A}$, a path of integration in $\bar{A}$ is projected to a path in $\underline{A}$.
 
First note that $m\ge 1$ since $\hat{A}$ contains the algebra $\{\alpha\mathbf{1}:\alpha\in\mathbb{R}\}$.  We cannot use $\mathcal{I}^T\mathcal{T}\,\mathcal{I}$  as employed in Theorem \ref{5/13/22.1} since in our present setting $\mathcal{T}$ does not necessarily map $A$ to $A$. However, using essentially the same argument as employed in the proof of Theorem \ref{5/13/22.1}, it is instead seen that $\Omega^T\,\mathcal{I}^T\mathcal{T}\,\mathcal{I}\,\Omega$ is an uncurling metric for the algebra $\bar{A}$.  Furthermore, analogous to the derivation of (\ref{9/27/23.1}) via (\ref{7/15/22.1}) and (\ref{8/25/22.1}), we obtain
 \begin{equation}\label{2/11/24.2}
 (\bar{s})^T \left[ \frac{\|\mathbf{1}_{\bar{A}}\|^2}{m} (\Omega^T\mathcal{I}^T\mathcal{T}\,\mathcal{I}\,\Omega)\right](\bar{s})^{-1}
  = \|\mathbf{1}_{\bar{A}}\|^2,\end{equation} so that the operator in brackets on the left-hand-side above is a normalized uncurling metric with respect to $\bar{A}$.  Taking $L$ in (\ref{11/7/20.2}) to be the bracketed expression in (\ref{2/11/24.2}), where $\mathbf{1}$ in (\ref{11/7/20.2}) is in this case $\mathbf{1}_{\bar{A}}$, equation (\ref{2/11/24.1}) follows analogously to (\ref{9/27/23.2}) in the proof of Theorem \ref{5/13/22.1}.
 \end{proof}
 
 The following is an immediate consequence.
 
 \begin{corollary}\label{1/1/23.1} An algebra's uncurling metric group is at least one-dimensional. \end{corollary}
  
Half of the examples of associative algebras $\bar{A}$ collected in Table I of Section 5 are such that the transpose operation applied to their left regular representation $A$ does not map to $A$, thereby illustrating Corollary \ref{11/17/22.1}.  Such calculations are explicitly highlighted in the examples of Section 4.4 and Section 4.5.3.

\section{Higher dimensional cases; application to some not associative algebras}
\subsection{Upper triangular Toeplitz matrix algebras}

The uncurling metric group and unital norm group associated with an upper triangular Toeplitz matrix algebra have some interesting features.
\begin{itemize}\item The uncurling metric group is given by the left regular representation (the algebra whose elements are lower triangular Toeplitz matrices) multiplied on the left by the exchange matrix. \item The usual norm is sensitive to only a single component of an element, but almost all unital norms are sensitive to all components of an element.\end{itemize}

\bigskip
 
 \begin{definition}\label{5/4/22.1} $\mathbb{T}_n$ is the $n$-dimensional vector space of real upper triangular Toeplitz $(n\times n)$-matrices. For $s\in \mathbb{T}_n$, the notation $s=(x_1,\dots,x_n)$ indicates that $s$ corresponds to the Toeplitz upper triangular $(n\times n)$-matrix whose entries on the successive diagonals, beginning with the main diagonal and proceeding to the right, are given by the successive vector component values $x_1,\dots,x_n$.   
\end{definition}

In the sequel, for $s\in \mathbb{T}_n$, $s$ can refer either to its format as an $n$-vector or its format as an $(n\times n)$-matrix, depending on the context.

The following result is known, but we include the proof for completeness, and to introduce some notation.  
 \begin{proposition}\label{5/10/22.1} Under the usual matrix product, $\mathbb{T}_n$ is a real $n$-dimensional unital associative algebra.\end{proposition} 
 
 \begin{proof}
Let ${\bf z}$ be the $(n \times n)$-matrix with entries having value $1$ on its superdiagonal  (the diagonal immediately to the right of the main diagonal) and value $0$ otherwise.  It is easily shown that for $j<n$ its $j$-th power, ${\bf z}^j$, is the $(n\times n)$-matrix with entries $1$ on the $j$-th diagonal to the right of the the main diagonal, with all other entries being $0$.   Thus, we can express the matrix format of $s\in\mathbb{T}_n$ as a polynomial in ${\bf z}$, such that
\begin{equation}\label{5/11/22.5} s = x_1 {\bf 1} + x_2 {\bf z} + x_3 {\bf z}^2 + \cdots + x_n {\bf z}^{n-1}.\end{equation}  It is similarly easy to show that ${\bf z}^k$ is the matrix of all zero entries if $k\ge n$.  It follows that for $s,t\in \mathbb{T}_n$, we have $st\in \mathbb{T}_n$.   Finally, for $x_1\ne 0$, we have \begin{equation}\label{5/11/22.6} s^{-1} =  (x_1 {\bf 1} + x_2 {\bf z} + x_3 {\bf z}^2 + \cdots + x_n {\bf z}^{n-1})^{-1}.\end{equation}  The matrix adding to $x_1{\bf 1}$ inside the parentheses on the right-hand-side of (\ref{5/11/22.6}) has only zero as an eigenvalue, so the full expression on the right-hand-side of (\ref{5/11/22.6}) has a power series expansion in ${\bf z}$.  Since ${\bf z}^k$ is the zero matrix for $k\ge n$, the power series expansion is  a $(n-1)$-degree polynomial in ${\bf z}$, which indicates that $s^{-1}\in \mathbb{T}_n$.
 \end{proof}
 
We will show that upper triangular Toeplitz matrix algebras have unital norms sensitive to all element components, providing a great contrast with the usual norm for these algebras, which is sensitive to only one element component.  One can easily generate other algebras with unital norms ``similarly dominating" the usual norm, as seen later.

To approach the problem of deriving the uncurling metric group for $\mathbb{T}_n$, the following will be important.

\begin{lemma} \label{5/7/22.15} For a unit $s\in \mathbb{T}_n$,  $n>1$, let $\left( s^{-1} \right)_p$ be the $p$-th component of $s^{-1}$ in the latter's $n$-vector format.  Given $k\le n$, then for all $j<k$ we have \begin{equation}\label{5/6/22.10} \frac{\partial \left( s^{-1} \right)_k}{\partial x_j} = \frac{\partial \left( s^{-1} \right)_{k-j+1}}{\partial x_1}.\end{equation}
\end{lemma}

\begin{proof} The proof is by induction. Equation (\ref{5/6/22.10}) is a tautology for the $k=2$ case.  We thus assume the inductive hypothesis that (\ref{5/6/22.10}) is true  for all cases up through $k-1$, and then establish its truth for the $k$-th case.

We must first express $s^{-1}$ in terms of the components of  $s=(x_1,\dots,x_n)\in \mathbb{T}_n$ with $x_1\ne 0$.   From the proof of Proposition \ref{5/10/22.1},  $s^{-1}$ is an upper triangular Toeplitz matrix.  But as an $n$-vector, it is easily seen that we can write it as  \begin{equation}\label{5/9/22.1} s^{-1} = (\frac{1}{x_1}, \frac{y_1}{x_1},\dots, \frac{y_{n-1}}{x_1}),\end{equation} where the $y_i$ are to be determined. 
Based on the notation introduced in the proof of Proposition \ref{5/10/22.1}, we can use (\ref{5/11/22.5}) and (\ref{5/11/22.6}) to write 
\begin{align}\nonumber s &= x_1 \left({\bf 1} + \frac{x_2}{x_1} {\bf z} + \frac{x_3}{x_1} {\bf z}^2 + \cdots + \frac{x_n}{x_1} {\bf z}^{n-1}\right)\\ s^{-1} &= x_1^{-1} \left({\bf 1} + y_1 {\bf z} + y_2 {\bf z}^2 + \cdots + y_{n-1} {\bf z}^{n-1}\right).  \label{5/18/22.10} \end{align}  Substituting the above two equations into $s^{-1} s = {\bf 1}$ leads to following set of relations, which allow the $y_i$ to be computed recursively:
\begin{align} -y_1 &= \frac{x_2}{x_1} \nonumber \\ -y_2 &= \frac{x_2}{x_1}y_1 + \frac{x_3}{x_1} \nonumber \\ &\;\vdots \nonumber \\ -y_{k-1} &= \frac{x_2}{x_1} y_{k-2} + \frac{x_3}{x_1} y_{k-3} + \cdots + \frac{x_k}{x_1} \label{5/7/22.1} \\ &\;\vdots \nonumber \end{align} 
 For $k>1$, (\ref{5/18/22.10}) implies  $\left( s^{-1} \right)_k = \frac{y_{k-1}}{x_1}$.  Applying this to (\ref{5/7/22.1}) with $k>1$, we find
\begin{eqnarray} \nonumber \left( s^{-1} \right)_k &=& -\frac{x_2}{x_1}\left( s^{-1} \right)_{k-1} - \frac{x_3}{x_1}\left( s^{-1} \right)_{k-2} -\dots  - \frac{x_k}{x_1}\left( s^{-1} \right)_1 \\ &=& -\frac{1}{x_1}\sum_{m=1}^{k-1} x_{m+1}\left( s^{-1} \right)_{k-m}. \label{5/7/22.2}\end{eqnarray}
Differentiating with respect to $x_j$, we find \begin{equation}  \frac{\partial \left( s^{-1} \right)_k}{\partial x_j} = -\frac{1}{x_1} \sum_{m=1}^{k-j} x_{m+1} \frac{\partial \left( s^{-1} \right)_{k-m}}{\partial x_j} -  \frac{1}{x_1}\left( s^{-1} \right)_{k-j+1} \label{5/7/22.4}, \end{equation} where the upper limit on the sum on the right-hand-side is $(k-j)$ rather than $(k-1)$ because the partial derivatives in the sum are zero for $m> k-j$ (i.e., $\left( s^{-1} \right)_p$ is not a function of $x_j$ for $p<j$).  

Next, substituting $(k-j+1)$ for $k$ in (\ref{5/7/22.2}) and then differentiating with respect to $x_1$, we have   \begin{eqnarray} \nonumber
 \frac{\partial \left( s^{-1} \right)_{k-j+1}}{\partial x_1} &=& \frac{1}{x_1^2} \sum_{m=1}^{k-j} x_{m+1} \left( s^{-1} \right)_{k-j-m+1}    - \frac{1}{x_1} \sum_{m=1}^{k-j} x_{m+1} \frac{\partial \left( s^{-1} \right)_{k-j-m+1}}{\partial x_1}  \\ &=& - \frac{1}{x_1}\left( s^{-1} \right)_{k-j+1}  -  \frac{1}{x_1} \sum_{m=1}^{k-j} x_{m+1} \frac{\partial \left( s^{-1} \right)_{k-j-m+1}}{\partial x_1},\label{5/7/22.3}\end{eqnarray} where the equality of the first terms on the right-hand-sides of the two above equalities comes from (\ref{5/7/22.2}) itself (with $(k-j+1)$ substituting for $k$). 
 Now, according to the induction hypothesis, $\frac{\partial \left( s^{-1} \right)_{k-m}}{\partial x_j} = \frac{\partial \left( s^{-1} \right)_{k-j-m+1}}{\partial x_1}$.  Inserting this into (\ref{5/7/22.3}), we see that the left-hand-sides of (\ref{5/7/22.4}) and (\ref{5/7/22.3}) are equal, as was required to be shown.
 \end{proof} 

Let $\slashed{\bf 1}_n$ be the exchange matrix, i.e., the $(n\times n)$-matrix with entries having value $1$ on the anti-diagonal and zero everywhere else.

\begin{theorem} \label{5/4/22.2} For any $t\in \mathbb{T}_n$ let $\hat{t}$ be its $(n\times n)$-matrix format..    Then $\hat{t} \slashed{\bf 1}_n$ is an uncurling metric.  Conversely, each member of the uncurling metric group associated with $\mathbb{T}_n$ can be expressed as $\hat{t}\slashed{\bf 1}_n$ for some $t\in \mathbb{T}_n$.   \end{theorem}

\begin{proof}  For $(\gamma_n,\dots,\gamma_1)\in \mathbb{T}_n$, let $G_n$ be its Toeplitz $(n\times n)$-matrix format (which happens to be the transpose of its left regular representation), and consider the resulting Hankel matrix, 
\begin{equation}\label{5/7/22.20} H_n \equiv G_n \slashed{\bf 1}_n = \begin{bmatrix} \gamma_1 & \gamma_2 & \gamma_3 & \dots & \gamma_n \\ \gamma_2 & \gamma_3 & \dots & \gamma_n & 0 \\ \gamma_3 & \dots & \iddots & \iddots & 0  \\ \vdots & \iddots & \iddots & \iddots & \vdots \\ \gamma_n & 0 & \dots & \ldots & 0 \end{bmatrix}.\end{equation}
For $\left( s^{-1} \right)_p$ as defined in the statement of Lemma \ref{5/7/22.15}, we then have
\begin{multline}  H_n s^{-1}   =  \Big( \gamma_1\left( s^{-1} \right)_1  + \dots + \gamma_n \left( s^{-1} \right)_n, \\ \gamma_2\left( s^{-1} \right)_1  + \dots + \gamma_n \left( s^{-1} \right)_{n-1}, 
  \dots , 
 \gamma_n \left( s^{-1} \right)_1 \Big). \label{5/7/22.40}
\end{multline}
 We will first demonstrate that $H_n$ is an uncurling metric.  According to (\ref{4/7/23.1}) as referred to by Definition \ref{9/20/23.1}, we are required to demonstrate \enlargethispage{.8cm}
\begin{equation}\label{5/7/22.30} d\left([H_n s^{-1}]\cdot ds\right) = 0.\end{equation}  

To verify (\ref{5/7/22.30}), we apply an induction argument.  The $n=2$ case is easily demonstrated to satisfy the above equation (e.g., see Section 4.4).  We assume the induction hypothesis that (\ref{5/7/22.30}) holds for all cases less than $n$, and now need to show that it holds for the $n$-th case.  
We begin by showing that the partial derivative of the first component of $H_n s^{-1}$ with respect to any $x_j$, with $j>1$, minus the partial derivative of the $j$-th component of $H_n s^{-1}$ with respect to $x_1$, is zero. We will then show that the analogous difference of partial derivatives of {\it any} two different components of $H_ns^{-1}$ is also zero.

So, from (\ref{5/7/22.40}), 
\begin{equation}\label{5/18/22.5} \frac{\partial \left(H_n s^{-1}\right)_j}{\partial x_1} = 
\sum_{i=j}^n \gamma_i\frac{\partial \left(s^{-1}\right)_{i-j+1}}{\partial x_1}. 
\end{equation}
But, \begin{equation}\label{5/18/22.6}  \frac{\partial \left(H_n s^{-1}\right)_1}{\partial x_j} = \sum_{i=1}^n \gamma_i\frac{\partial \left(s^{-1}\right)_i}{\partial x_j}
= \sum_{i=j}^n \gamma_i\frac{\partial \left(s^{-1}\right)_i}{\partial x_j} = \sum_{i=j}^n \gamma_i\frac{\partial \left(s^{-1}\right)_{i-j+1}}{\partial x_1},
\end{equation} where the second equality follows (\ref{5/9/22.1}) and (\ref{5/7/22.1}), which indicate that $\frac{\partial \left(s^{-1}\right)_i}{\partial x_j} = 0$ for $i<j$, and the third equality follows from Lemma \ref{5/7/22.15}.
Equations (\ref{5/18/22.5}) and (\ref{5/18/22.6})  demonstrate that the partial derivative of the first component of $H_n s^{-1}$ with respect to any $x_j$, with $j>1$, minus the partial derivative of the $j$-th component of $H_n s^{-1}$ with respect to $x_1$, is equal to zero.

Now we will consider the vanishing of the partial derivative of the $k$-th component of $H_n s^{-1}$ with respect to $x_j$ minus the partial derivative of the $j$-th component of $H_n s^{-1}$ with respect to $x_k$, for all possible $k>j>1$ (i.e, all choices relevant to (\ref{5/7/22.30}) other than considered in the prior paragraph).  In fact, based on the telescoping format of (\ref{5/7/22.40}), we already have the result that this difference vanishes, because it merely reduces to an instance of the theorem for the case of a matrix of the form $H_ks^{-1}$ with $k<n$ - and this case is assumed to hold according to the induction hypothesis.  So, (\ref{5/7/22.30}) is established.

Furthermore, a Hankel matrix of the form (\ref{5/7/22.20}) is the only possibility for an uncurling metric $L$ because any other choice resulting in the coefficients of the terms in any two components of $Ls^{-1}$ cannot lead to a null value for the difference of the two requisite partial derivatives (as is necessary for $d\left([L s^{-1}]\cdot ds\right) = 0$) due to linear independence of the individual terms summing to the components. That is, replacing  $H_n$ with matrix $L=(a_{p,q})_{p,q=1}^n$ in (\ref{5/18/22.5}) and (\ref{5/18/22.6}), and invoking Lemma \ref{5/7/22.15}, we obtain
\begin{eqnarray}\nonumber \frac{\partial \left(L s^{-1}\right)_j}{\partial x_1}  -  \frac{\partial \left(L s^{-1}\right)_1}{\partial x_j} &=&  \sum_{i=1}^n a_{j,i}\frac{\partial \left(s^{-1}\right)_i}{\partial x_1} - \sum_{i=j}^n a_{1,i}\frac{\partial \left(s^{-1}\right)_i}{\partial x_j} \\  &=& \sum_{i=1}^n a_{j,i}\frac{\partial \left(s^{-1}\right)_i}{\partial x_1} - \sum_{i=j}^n a_{1,i}\frac{\partial \left(s^{-1}\right)_{i-j+1}}{\partial x_1} \nonumber \\
&=& \sum_{i=1}^n a_{j,i}\frac{\partial \left(s^{-1}\right)_i}{\partial x_1} - \sum_{i=1}^{n-j+1} a_{1,i+j-1}\frac{\partial \left(s^{-1}\right)_i}{\partial x_1}.
\label{5/9/22.10}
\end{eqnarray}
This indicates that we must have $a_{j,i}=0$ for $i>n-j+1$, and otherwise have $a_{j,i} = a_{1,i+j-1}$.  Since this must hold for any $j>1$, $L$ must be a Hankel matrix of the form (\ref{5/7/22.20}).  Thus, the $H_n$ are the only possibilities for uncurling metrics. 

Since $t=(\gamma_n,\dots,\gamma_1)$ is an arbitrary member of $\mathbb{T}_n$, the theorem follows.
\end{proof}

\begin{corollary}\label{7/1/22.1} For any $t\in \mathbb{T}_n$, denote the value of the $(1,n)$ entry of its matrix format $\hat{t}$ as $\hat{t}_{1,n}$.  Then $\frac{1}{t_{1,n}}\hat{t}\slashed{\bf 1}_n$ is a normalized uncurling metric if $\hat{t}_{1,n}\ne 0$, and all normalized uncurling metrics can be written in this form for some $t\in \mathbb{T}_n$. \end{corollary}
\begin{proof} Theorem \ref{5/4/22.2} has shown that all uncurling metrics of $\mathbb{T}_n$ have the form (\ref{5/7/22.20}).
We now consider constraint (\ref{7/7/21.1}) of Definition \ref{9/20/23.1}.  Let $H_n^{(1)}$ be matrix $H_n$ except with the  $(1,1)$-entry having the value $1$.
Since $\|{\bf 1}\|^2 =1$ for this algebra, we must show that \[ s^T H_n^{(1)} s^{-1} = s\cdot [H_n^{(1)}s^{-1}] = 1.\]  From (\ref{5/9/22.1}) and (\ref{5/7/22.20}), we have
\begin{eqnarray}\nonumber s^T H_n s^{-1} &=& (x_1,\dots,x_n)\cdot \left[ H_n \left(\frac{1}{x_1}, \frac{y_1}{x_1},\dots,\frac{y_{n-1}}{x_1}\right)^T\right] \\ 
 &=&  \gamma_1 + \gamma_2y_1 + \gamma_3y_2 \hspace{3mm} + \hspace{3mm}\gamma_4y_3 \hspace{3mm} + \dots   \hspace{6mm}+\gamma_ny_{n-1} \nonumber \\
&+&  \gamma_2 \frac{x_2}{x_1} + \gamma_3\frac{x_2}{x_1}y_1 + \hspace{2mm} \gamma_4\frac{x_2}{x_1}y_2 + \dots +  \gamma_n\frac{x_2}{x_1}y_{n-2} \nonumber \\ &+&\gamma_3 \frac{x_3}{x_1}  + \hspace{3mm}\gamma_4\frac{x_3}{x_1}y_1 +  \dots \hspace{4mm} +  \gamma_n\frac{x_3}{x_1}y_{n-3} \nonumber \\ 
&+&  \hspace{5mm} \vdots \hspace{2.5cm} \iddots \bigskip\bigskip \nonumber \\ &+& \gamma_{n-1} \frac{x_2}{x_1} \hspace{3mm} + \gamma_n \frac{x_2}{x_1}y_1 \nonumber \\ &+& \hspace{2mm} \gamma_n\frac{x_n}{x_1}. \label{6/23/22.1}
 \end{eqnarray}
 For any $i>1$, if we collect all the terms on the right-hand-side above that each have $\gamma_i$ as a coefficient, and then apply (\ref{5/7/22.1}), it is quickly seen that the sum of all terms having $\gamma_i$ as a coefficient vanishes.  It follows that
 $s^T H_n s^{-1} =\gamma_1$.   Thus, the requisite matrix satisfying constraint (\ref{7/7/21.1}) of Definition \ref{9/20/23.1} will be the right-hand-side of (\ref{5/7/22.20}) with $\gamma_1=1$.  Setting $\frac{1}{t_{1,n}}t = (\gamma_n,\dots,\gamma_2, 1)$, the corollary follows.
\end{proof}

\begin{corollary}\label{5/17/22.1} For $s= (x_1,\dots,x_n)\in \mathbb{T}_n$ with $x_1\ne 0$, the usual norm is $x_1^n$.  The special unital norm family of $\mathbb{T}_n$ is the $(n-1)$-parameter family,
\begin{equation} \label{5/5/22.6} \ell_{\gamma_2,\dots,\gamma_n}(x_1,\dots,x_n) = x_1 e^{P_{\gamma_2,\dots,\gamma_n}\left(\frac{x_2}{x_1},\dots,\frac{x_n}{x_1}\right)},\end{equation} \enlargethispage{1cm} where $P_{\gamma_2,\dots,\gamma_n}(w_1,\dots,w_{n-1})$ is a real $(n-1)$-degree multivariate polynomial with coefficients dependent on the $(n-1)$ real parameters $\gamma_2,\dots,\gamma_n$.    \end{corollary}

\begin{proof} The usual norm of $s=(x_1,\dots,x_n)\in \mathbb{T}_n$ is immediately seen to be $x_1^n$.

With respect to the notation $G_n,H_n$ introduced at the beginning of the proof of Theorem \ref{5/4/22.2}, each $(\gamma_n,\dots,\gamma_2,1)\in\mathbb{T}_n$ is associated with a Toeplitz matrix $G_n^{(1)}$ having its $(1,n)$-entry equal to $1$.  Corollary \ref{7/1/22.1} indicates that $H^{(1)}_n\equiv G_n^{(1)}\slashed{\bf 1}_n$ is a normalized uncurling metric. From (\ref{11/7/20.2}),  the associated unital norm is $e^{\left[\int_{\bf 1}^s (H_n^{(1)} v^{-1})\cdot dv\right]}$.  Elementary integration and a simple induction argument quickly results in (\ref{5/5/22.6}) and the $(n-1)$-degree of the multivariate polynomial in question.
\end{proof}
 
\subsection{Extension to certain not associative algebras}
The two-dimensional real unital algebras have generalizations as the Cayley-Dixon sequence algebras ($\mathbb{R}$ and $\mathbb{C}$ being its lowest dimensional members),  the Spin Factor Jordan algebras (generalizing $\slashed{\mathbb{C}}$), and real upper triangular matrix algebras with at least some redundant diagonal components (generalizing the lowest dimensional case, $\mathbb{D}$).  Their unital norm groups will be the subject of this subsection. 

Spin Factor Jordan Algebras and Cayley-Dixon Algebras after the quaternions are power-associative but not associative.  While the usual norm is thereby inapplicable, there is an ``Algebraic Norm" replacement \cite{jacobson:1963} (based on the minimal polynomial, and thus analogous to the usual norm as given by the constant term of a characteristic polynomial).  For the Spin Factor Jordan Algebras and Cayley-Dixon Algebras, the Algebraic Norms are essentially the  Minkowski  norm and Euclidean norm, respectively.  Thus, useful norms are not restricted to associative (or alternative) algebras, even though it turns out that the multiplicative property (i.e., $N(s_1s_2)=N(s_1)N(s_2)$) must then be sacrificed  \cite{schafer:1963}. 
So, despite a strong bias against a norm not satisfying the multiplicative property, these examples indicate that ``algebraic" norms that are not multiplicative can nevertheless be very significant \cite{mccrimmon:2004}. 

Of course, apart from the omnipresent unital norm group member comparable to the usual norm as guaranteed by Corollary \ref{11/17/22.1} for real finite-dimensional unital associative algebras (as well as all unital norms resulting from two-dimensional real unital algebras), unital norms are not typically multiplicative.  The multiplicative property is instead replaced by (\ref{5/23/23.3}).  The latter equation also replaces the multiplicative property for the Algebraic Norm in the cases of the not associative Spin Factor Jordan Algebras and Cayley-Dixon Algebras. 
Like the Algebraic Norm, the unital norm/uncurling metric constructions remain valid for these latter two classes of algebras.  The way we deal with these particular not associative algebras is as follows.

\subsubsection{Unital norms for $^*$-algebras whose Hermitian elements are real} 

Consider an algebra associated with an involutive antiautomorphism ``
$^*$ " (known as a $^*$-algebra, star-algebra, or involutive algebra) such that the Hermitian elements are members of $\mathbb{R}$ (thus, $s+s^*$ and $s^*s$ are both real; also, $s^*s = ss^*$ because $(s+s^*)s = s(s+s^*)$ since $(s+s^*)$ is real). 
Then \[ X^2 - (s+s^*)X + s^*s, \] is the minimal polynomial of any not real element $s$, and annihilates all elements, real and not real.  Consistent with comments in the second paragraph of Section 4.2, we define its constant term to be the ``Algebraic Norm" of an element $s$, by analogy with the characteristic polynomial defined for an associative algebra, which supplies the norm from the constant term of that element-associated polynomial annihilating all elements.


On these $^*$-algebras, the multiplicative identity $\mathbf{1}$ is $1$ (where the latter is thought of as the algebra element that is the multiplicative identity on the subalgebra $\mathbb{R}$ rather than a field element).  Accordingly, for these algebras, we define the expression $\|{\bf 1}\|^2$ (e.g., that appears in Definition \ref{9/20/23.1}) to be $1$.  We also have,

\begin{lemma}\label{9/14/22.2} The inverse of a unit $s$ in a finite-dimensional $^*$-algebra whose Hermitian elements are real is unique, and given by $s^{-1}=\frac{s^*}{ss^*}$.\end{lemma}

\begin{proof} 
We can write $s=\frac{s+s^*}{2} +\frac{s-s^*}{2}$, where the first  term on the right-hand-side is real and the second term on the right-hand-side is antisymmetric with respect to application of the $^*$-operation.  We define $\sigma$ to be the first term on the right-hand-side, and $\bar{s}$ to be the second term on the right-hand-side.  Suppose $x = \chi + \bar{x}$ is an inverse of $s$, where $\chi$ is real, and $\bar{x}$ is antisymmetric. Then \begin{equation}\label{9/14/22.1} sx = \sigma\chi + \sigma\bar{x}+\chi\bar{s} + \bar{s}\bar{x} = 1,\end{equation} and $xs = \sigma\chi + \sigma\bar{x}+\chi\bar{s} + \bar{x}\bar{s} = 1$.  Thus, $\bar{s}\bar{x} = \bar{x}\bar{s}$.  Since $\bar{s}^* = (\frac{s-s^*}{2})^* = -(\frac{s-s^*}{2}) = -\bar{s}$, we have $(\bar{s}\bar{x})^* = \bar{x}^*\bar{s}^* = \bar{x}\bar{s} = \bar{s}\bar{x}$.  Hence, $\bar{s}\bar{x}$ is real.  It follows from (\ref{9/14/22.1}) that $\sigma\chi+\bar{s}\bar{x} =1$ and $\sigma\bar{x}+\chi\bar{s} = 0$.  The latter has has a unique solution, yielding $\chi+\bar{x} = \frac{\sigma-\bar{s}}{\sigma^2-\bar{s}\bar{s}} = \frac{s^*}{ss^*}$, where the second equality follows since the right-hand-side is obviously an inverse ($1=s\frac{s^*}{ss^*} = \frac{s^*}{ss^*}s$ because $s^*s=ss^*$ is real and thus self-adjoint), and the inverse has just now been seen to be unique.
\end{proof}

Extending Definition \ref{9/20/23.1} to include this class of $^*$-algebras, we can now provide a special unital norm  by identifying a linear transformation $L$ satisfying (\ref{4/7/23.1}), (\ref{7/7/21.1}) and then evaluating (\ref{11/7/20.2}).  
 
 \begin{proposition}\label{1/11/22.1} Given a $^*$-algebra with vector space of elements $\mathbb{R}^n$ and whose Hermitian elements are members of $\mathbb{R}$, the Algebraic Norm on the $^*$-algebra in a sufficiently small neighborhood of ${\bf 1}$ is the square of a special unital norm. \end{proposition}

\begin{proof} By definition, the Algebraic Norm is \begin{equation}\label{1/15/22.1} N(s)= ss^* = N(s^*)\in\mathbb{R}.\end{equation} 
 We then have $N(\alpha s) = \alpha^2 N(s)$ for $\alpha\in \mathbb{R}$, i.e., degree-2 homogeneity of $N(s)$.  
For $N(s)\ne 0$, Lemma \ref{9/14/22.2} implies 
\begin{equation}\label{1/12/22.1} s^{-1} = \frac{s^*}{N(s)}.\end{equation}
Applying $N(\cdot)$ to both sides of  (\ref{1/12/22.1}), the degree-2 homogeneity and (\ref{1/15/22.1}) imply $N(s)N(s^{-1})=1$, so that \begin{equation}\label{1/15/22.2} \sqrt{N(s)}\sqrt{N(s^{-1})}=1. \end{equation}

We are required to show that $\sqrt{N(s)}$ is a special unital norm.  Thus, we must first verify that $\sqrt{N(s)}$ satisfies the requirements of the extended  version of Definition \ref{9/20/23.1} as indicated prior to the statement of this proposition. 

To begin with, (\ref{1/15/22.2}) already satisfies  (\ref{5/23/23.3}).
 Equation (\ref{1/15/22.1}) indicates that $N(s)$ is a quadratic form, so there is an orthonormal basis with respect to which the members of the algebra can be written as $s= (s_1,\dots, s_i,\dots)$ with $N(s) = \sum_i \delta_i s_i^2$ and $\delta_i = -1$, $0$, or $1$, since $^*$ is an involution.  Let $K_1 \equiv {\rm diag}\{\delta_i\}$.   Then, \begin{equation}\label{2/6/22.1} N(s) = s\cdot(K_1 s).\end{equation}
If $s^{-1}$ exists (i.e., $N(s)\ne 0$), we then have \begin{equation}\label{2/6/22.2} \nabla\left[\sqrt{N(s)}\right] = \frac{K_1s}{\sqrt{N(s)}}.\end{equation}  Define $K_2$ to be the algebra's involutive antiautomorphism operation $^*$.  It follows that $K_1K_2K_2[\frac{s}{\sqrt{N(s)}}] = \nabla\left[\sqrt{N(s)}\right]$. Defining $L\equiv K_1K_2$, dividing both sides of the equation in the last sentence by $\sqrt{N(s)}$, and using (\ref{1/15/22.1}) and (\ref{1/12/22.1}), we then have \begin{equation}\label{1/12/22.2} Ls^{-1}  = \frac{\nabla\left[\sqrt{N(s)}\right]}{\sqrt{N(s)}} .\end{equation} This equation  is valid for $N(s)>0$, and the latter pertains in a sufficiently small neighborhood of ${\bf 1}$.  Thus,  (\ref{11/7/20.1}) is satisfied.  From (\ref{2/6/22.1}), (\ref{2/6/22.2}), and (\ref{1/12/22.2}), we have \[ s\cdot [Ls^{-1}] = s\cdot K_1\left[\frac{s}{N(s)}\right] = 1,\] which is (\ref{7/7/21.1}).  
 
 To establish that $\sqrt{N(s)}$ is a unital norm, we  now only need to verify that (\ref{4/7/23.1}) is satisfied, i.e., $d\left([Ls^{-1}]\cdot ds\right)=0$.  According to (\ref{2/6/22.2}), (\ref{1/12/22.2}), and the definition  of $K_1$, we have $d\left([Ls^{-1}]\cdot ds\right) = d\left(\sum_i \frac{\delta_is_i ds_i}{N(s)}\right)$.  Performing the requisite differentiation to obtain the component of the exterior derivative associated with a choice of $i$ and $j$, we find \[ \frac{\partial [Ls^{-1}]_j}{\partial s_i} - \frac{\partial [Ls^{-1}]_i}{\partial s_j} = -2\frac{(\delta_j s_j)(\delta_i s_i)}{N^2(s)} +
2\frac{(\delta_i s_i)(\delta_j s_j)}{N^2(s)} = 0,  \] where $\frac{\partial [Ls^{-1}]_j}{\partial s_i}$ is the partial derivative of the $j$-th component of $[Ls^{-1}]$ with respect to the $i$-th component of $s$, etc. Thus,  $d\left([Ls^{-1}]\cdot ds\right) = 0$.

Consequently, $\sqrt{N(s)}$ satisfies the special unital norm requirements of the extended version of Definition \ref{9/20/23.1}.  
\end{proof}


\subsubsection{The Cayley-Dixon sequence}  From the proof of Proposition \ref{1/11/22.1}, it is easily shown that a normalized uncurling metric associated with an algebra of the Cayley-Dixon sequence (all of whose members are $^*$-algebras whose Hermitian elements are real) is given by ${\rm diag}\{1,-1,-1,\dots,-1\}$, with the associated special unital norm seen to be the Euclidean norm. Although it is easily shown from results in Section 2.3 that $\mathbb{C}$ is associated with a 1-parameter family of normalized uncurling metrics, a simple argument (using Lemma \ref{9/14/22.2}) indicates that all other algebras of the sequence have only the above normalized uncurling metric, i.e., the special unital norm family has only one member.  
 
\subsubsection{Spin Factor Jordan Algebras}  Consider the bilinear product ``$\,\bullet\,$" such that for
$\alpha,\beta\in \mathbb{R},\,\, {\sf a},{\sf b}\in \mathbb{R}^m$,  \begin{equation}  (\alpha+ {\sf a})\bullet (\beta + {\sf b}) \equiv
    [\alpha \beta + {\sf a} \cdot {\sf b}] + [ \beta {\sf a} + \alpha {\sf b}] \in \mathbb{R}\oplus \mathbb{R}^m. \label{12/24/18.1}\end{equation}  This makes the vector space $\mathbb{R}\oplus \mathbb{R}^m$ a unital algebra with the multiplicative identity ${\bf 1}$ as simply $1$ (where in this instance $1$ is understood to be an element of the algebra rather than a member of the field).  For $m\ge 2$, this is known as a Spin Factor Jordan Algebra \cite{mccrimmon:2004}.

 There are two involutive antiautomorphisms on each of these commutative algebras, one of which is the identity and the other is the conjugate operation, $(\alpha+ {\sf a})^* = (\alpha- {\sf a})$. The latter is the only one of the two involutive antiautomorphisms  for which the Hermitian elements are real.  Applying Proposition \ref{1/11/22.1}, it is easily shown that the Algebraic Norm is the special unital norm given by the Minkowski norm.  Similar to the case of the Cayley-Dixon algebras, it is easy to show that this is the only special unital norm.

 
 \subsection{The algebra with vector space of elements $\mathbb{R}^n$ and component-wise multiplication} A higher dimensional analogue of $\mathbb{R}\oplus \mathbb{R}$, taken as the algebraic direct sum of $n$ copies of $\mathbb{R}$, is associated with an $(n-1)$-parameter family of  special unital norms.

 This algebra is isomorphic to the subalgebra of $ M_n(\mathbb{R})$  whose members are diagonal matrices. In this example, ${\bf 1} = (1,\dots,1)$ and $\|{\bf 1}\|^2 = n$.   It is easily appreciated that (\ref{4/7/23.1}) forwards the uncurling metric group ${\rm diag} \{ \sigma_1, \sigma_2,\dots,\sigma_n\}$, so the uncurling metric group is the same as the members of the left regular representation of the algebra.  This means that for $s=(x_1,\dots,x_n)$, we must have $L[(x_1,x_2,\dots,x_n)^{-1}] = (\frac{\sigma_1}{x_1}, \frac{\sigma_2}{x_2},\dots,\frac{\sigma_n}{x_n})$.     Evaluation of (\ref{11/7/20.2}) indicates that the members of the unital norm group are given by  \begin{equation}\label{1/15/22.5} \ell_{\sigma_1,\dots,\sigma_n}(x_1,\dots,x_n) = (x_1^{\sigma_1} x_2^{\sigma_2} \dots x_n^{\sigma_n})^\frac{1}{n}. \end{equation} From (\ref{7/7/21.1}), the special unital norms are given by (\ref{1/15/22.5}) under the constraint $\sum_1^n \sigma_i = n$. Thus, for this algebra we have a $(n-1)$-dimensional special unital norm family.   
 
 Similar to the analogous case in two dimensions (Section 2.3), members of the unital norm group exist that are sensitive only to any desired subset of an element's components, and all units can potentially be reconstructed from combinations of different unital norm values - unlike with the usual norm.
 
 The geometric mean is the special unital norm given by (\ref{1/15/22.5}) with $\sigma_1=\sigma_2=\dots =\sigma_n = 1$.  
   
 \subsection{Upper triangular matrix algebras with independent entries on the main diagonal} Consider the real three-dimensional algebra with elements $(x,y,z)$ which multiply according to the matrices 
$\left[
\begin{array}{ccc}
 x & z   \\
 0 & y    
\end{array}
\right]$.  Or in other words, the algebra is $\left\{ \left[
\begin{array}{ccc}
 x & z   \\
 0 & y    
\end{array}
\right],x,y,z\in\mathbb{R}\right\}$.  We have $(x,y,z)^{-1} = \left( \frac{1}{x}, \frac{1}{y}, -\frac{z}{xy} \right)$, and ${\bf 1}= (1,1,0)$, and $\|{\bf 1}\|^2 = 2$.  Starting with a generic symmetric $L = \left[
\begin{array}{ccc}
 \alpha & \beta & \gamma   \\
 \beta & \delta & \epsilon \\ \gamma & \epsilon & \xi    
\end{array}
\right]$, 
it is required that $\left( \frac{\epsilon(y-x)+\xi z}{xy^2}, \frac{\gamma(y-x)-\xi z}{x^2y}, \frac{\beta(x^2-y^2)-\gamma x z +\epsilon y z}{x^2y^2}\right)$ vanishes in a neighborhood of $(1,1,0)$.  This implies that all entries of the matrix $L$ must be zero except for $\alpha$ and $\delta$.   Thus, $L(x,y,z)^{-1} = (\frac{\alpha}{x}, \frac{\delta}{y},0)$.  From (\ref{7/7/21.1}), we have the requirement that a normalized uncurling metric must satisfy \[ (x,y,z)\cdot [L(x,y,z)^{-1}] = \alpha +\delta =  2.\]  Equivalently, we can alternatively specify $\alpha \equiv 1+\sigma$ and $\delta \equiv 1-\sigma$, for arbitrary $\sigma\in\mathbb{R}$.   We thus have the normalized uncurling metric family composed of all  \begin{equation}\label{11/17/22.2} L = \left[
\begin{array}{ccc}
 1+\sigma & 0 & 0   \\
 0 & 1-\sigma & 0 \\ 0 & 0 & 0    
\end{array}
\right],\end{equation} $\sigma\in\mathbb{R}$.  In this case, each member of the family is singular.  Equation (\ref{11/7/20.2}) then leads to the 1-parameter family of special unital norms, \begin{equation}\label{12/21/21.1} \ell_\sigma(x,y,z) = x^\frac{1+\sigma}{2} y^\frac{1-\sigma}{2} = (xy)^{\frac{1}{2}}\left(\frac{x}{y}\right)^{\frac{\sigma}{2}}.\end{equation} 

The left regular representation of an element $(x,y,z)$ is given by $\left[\begin{array}{ccc}
x & 0 & 0   \\
 0 & y & 0 \\ 0 & z & x    
\end{array} \right]$.  Thus, the $\mathbb{R}^{3\times 3}$ transpose operation maps outside the algebra's representation.  However, for the subalgebra defined by elements $(x,y,0)$, the left regular representation is 
 $\left[\begin{array}{ccc}
x & 0  \\
 0 & y  
\end{array} \right]$.  The usual norm of the subalgebra is thus $xy$, and one of the special unital norms on the subalgebra is $(xy)^{\frac{1}{2}}$, and these differ only by an exponent - consistent with Theorem \ref{5/13/22.1}.   Note that the latter unital norm on the subalgebra provides the values for a unital norm on the full algebra, corresponding to the choice of normalized uncurling metric in (\ref{11/17/22.2}) given by $\sigma = 0$, i.e., $\ell_0(x,y,z) = (xy)^{\frac{1}{2}}$, consistent with Corollary \ref{11/17/22.1}.  The fact that the transpose operation on $\mathbb{R}^{3\times 3}$ does not map the above left regular representation of $\left\{\left[\begin{array}{ccc}
x & z  \\
 0 & y  
\end{array} \right]:(x,y,z)\in \mathbb{R}^3\right\}$ to itself explains why there is no unital norm that is a power of the usual norm $x^2y$.
  
\subsection{Some higher dimensional algebras with unital norms that are sensitive to more element components than is  usual norm}

The usual norm on $ M_n(\mathbb{R})$ is sensitive to all components of an element of the algebra (because the determinant of an element of $ M_n(\mathbb{R})$  is sensitive to the values of all entries of its matrix).    However, this is not true in general for associative algebras.  For example, for any algebra isomorphic to a triangular matrix algebra, the usual norm is only sensitive to the components on the main diagonal of the triangular matrix representation.  This is also true for the unital norm if all entries on the main diagonal of the triangular matrix representation are allowed to be distinct.  But if  at least two entries on the main diagonal always take equivalent values for the particular algebra,  then almost all of its unital norms family members will be sensitive to one or more components not on the main diagonal, in addition to all of the components on the main diagonal (these are matrix algebras for which an element that is not a diagonal matrix is never diagonalizable).  

It is easily seen from the presentation in Section 2.3 that the dual numbers are a two-dimensional example of this, and it
generalizes to higher dimensions according to Corollary \ref{5/17/22.1}.  We will now highlight an additional feature using the next three examples.  

\subsubsection{A three-dimensional algebra}Consider the real algebra with elements  $(x,z,w)$ that multiply according to the matrices $ \left[
\begin{array}{ccc}
x & z & w   \\
0 & x & z \\
0 & 0 & x  
\end{array}
\right]$.    The constraints (\ref{7/21/23.1}) and (\ref{7/7/21.1}) lead to the normalized uncurling metrics, $ L = \left[
\begin{array}{ccc}
1 & \beta & \gamma   \\
\beta & \gamma & 0 \\
\gamma & 0 & 0  
\end{array}
\right], \,\,\beta,\gamma\in\mathbb{R}$, which is just the $n=3$ case of Corollary \ref{7/1/22.1}.    Evaluating (\ref{11/7/20.2}), we now have a 2-parameter family of special unital norms, \begin{equation}\label{3/31/21.3} \ell_{\beta,\gamma}(x,z,w) = x\, e^{ \beta\frac{z}{x}  +\gamma \frac{w}{x} -\gamma \frac{ 1}{2}\left(\frac{z}{x}\right)^2 }.\end{equation}    The choice $\beta = \gamma = 0$ produces the special unital norm, $\ell_{0,0}(x,z,w)= x$, which is comparable to the usual norm, $x^3$. 
 The other special unital norms are distinguished by their dependence on $z$ and/or $w$, i.e., unlike the usual norm.  The non-units are associated with a different type of unital norm singularity (as $x$ approaches 0) than seen in examples considered so far other than $\mathbb{T}_n$.
 
 \subsubsection{A four-dimensional algebra} Consider the real unital algebra with elements 
$(x,z,v,w)$ which multiply according to the matrices 
$\left[
\begin{array}{ccccc}
x & z & w   \\
0 & x & v \\
0 & 0 & x  
\end{array}
\right]$.
Straightforward application of the constraints (\ref{7/21/23.1}) and (\ref{7/7/21.1}) leads to the normalized uncurling metric family 
$\left[
\begin{array}{cccc}
1 & \beta & \gamma & \delta   \\
\beta & 0 & \frac{1}{2}\delta & 0 \\
\gamma & \frac{1}{2}\delta &  0 & 0   \\
\delta & 0 &  0 & 0
\end{array}
\right]$, and the associated special unital norm family,
\[ \ell_{\beta,\gamma,\delta}(x,v,z,w) = x\, e^{\beta\frac{v}{x} + \gamma\frac{z}{x} +\delta\frac{w}{x}- \frac{\delta}{2}\left(\frac{zv}{x^2}\right)}.\]


\subsubsection{A five-dimensional algebra}Consider the real unital algebra with elements 
$(x,y,z,v,w)$ which multiply according to the matrices
$\left[
\begin{array}{ccccc}
x & z & w   \\
0 & x & v \\
0 & 0 & y  
\end{array}
\right]$.
Straightforward application of the constraints (\ref{7/21/23.1}) and (\ref{7/7/21.1}) leads to the normalized uncurling metric family 
$\left[
\begin{array}{ccccc}
1+\sigma & 0 & \beta & 0  & 0   \\
0 & 1-\sigma & 0 & 0 & 0 \\
\beta & 0 &  0 & 0 & 0  \\
 0 & 0 &  0 & 0 & 0  \\
0 & 0 &  0 & 0 & 0
\end{array}
\right]$, and the associated special unital norm family 
\[ \ell_{\sigma,\beta}(x,y,z,v,w) =  (xy)^{\frac{1}{2}}\left(\frac{x}{y}\right)^{\frac{\sigma}{2}} e^{\frac{\beta}{2} \frac{z}{x}}.\]

In this example, almost all unital norms are influenced the element's $z$-component, while this sensitivity is not present with the usual norm (given by $x^4y$).  However, like the usual norm,  the unital norm remains insensitive to the $w$ and $v$ components.  Once again, there is a new unital norm singularity type as non-units are approached.  

Finally, the left regular representation of this algebra with elements $(x,y,z,v,w)$ is given by $ \left[
\begin{array}{ccccc}
x & 0 & 0 & 0 & 0   \\
0 & y & 0 & 0 & 0 \\
z & 0 &  x & 0 & 0  \\
0 & v &  0 & x & 0  \\
0 & w &  0 & z & x
\end{array}
\right].$
The transpose operation on $\mathbb{R}^{5\times 5}$ does not map this representation algebra into itself, due to its effect on the last three columns and last three rows - which explains why there is no unital norm that is a power of the usual norm $x^4y$.  However, modifying the transpose operation to act only on the $(2\times 2)$-submatrix of values in the upper left corner of the representation (while suppressing the components $v,z,w$) produces the unital norm $(xy)^\frac{1}{2}$, consistent with Corollary \ref{11/17/22.1}.
 
\section{Additional observations relevant to epimorphisms and classification}

\subsection{Epimorphisms} Table I collects the previously presented examples of uncurling metric groups and associated special unital norms.  Greek letters represent parameters ranging over $\mathbb{R}$.  The Spin Factor Jordan Algebras and Cayley-Dixon sequence algebras that are not associative are treated via their formulation as \mbox{$^*$-algebras} whose Hermitian elements are real, as in Section 4.2. The uncurling metrics and special unital norms are written with respect to the standard basis.  Each of the fourteen rows of Table I refers to the algebra given as the first entry of the row.  $\slashed{\mathbb{C}}$ is the algebra of split complex numbers, $\mathbb{D}$ is the algebra of the dual numbers, and $\mathbb{H}$ is the quaternion algebra.  $\mathbb{T}_n$ is the $n$-dimensional upper triangular Toeplitz matrix algebra.  The notation $A_q$ refers to the algebra in the $q$-th row.  In the third entry of row 14, the exponent is a multivariate polynomial in the indicated variables, with the subscripted parameters multiplying constituent monomial terms.  $\slashed{\mathbf{1}}_n$ is the $(n\times n)$-exchange matrix.  The expressions $\mathbb{C}_\text{left}$, $\slashed{\mathbb{C}}_\text{left}$, $[\mathbb{R}\oplus\mathbb{R}]_\text{left}$, $\mathbb{D}_\text{left}$, and $[\mathbb{T}_n]_\text{left}$, indicate the left regular representations.  


\bigskip

The examples in Table I illustrate many of the points made earlier in this paper.    
 
\begin{itemize}\item $\mathfrak{u}_{M_n(\mathbb{R})}$ is one-dimensional (Theorem 3.1), from which one concludes that $M_n(\mathbb{R})$ is simple (Corollary \ref{8/17/22.1}) - a demonstration distinct from the usual one.
\item Up to an isomorphism, there are only three two-dimensional real unital algebras, e.g., $\mathbb{C}$, $\slashed{\mathbb{C}}$, and $\mathbb{D}$.  Algebras $\mathbb{R}\oplus\mathbb{R}$ and $\slashed{\mathbb{C}}$ are isomorphic via the linear transformation $\left[\begin{array}{ccc}
 1 & -1   \\ 
 1 & 1    
\end{array}\right]$, but they are not isometrically isomorphic, since the isomorphism does not have unit determinant.  The respective left regular representations of $\mathbb{C}$, $\slashed{\mathbb{C}}$, and $\mathbb{D}$, are $\left[\begin{array}{ccc}
 x & -y   \\ 
 y & x    
\end{array}\right]$, $\left[\begin{array}{ccc}
 x & y   \\ 
 y & x    
\end{array}\right]$, and $\left[\begin{array}{ccc}
 x & 0   \\ 
 y & x    
\end{array}\right]$.  For each of the algebras $\mathbb{C}$ and $\mathbb{D}$, the set of members of the algebra's uncurling metric group is the same as the set of members of the algebra's left regular representation algebra except with each of the latter's matrix members multiplied on the left by the exchange matrix $\slashed{\mathbf{1}}_2$ (a similar statement holds for $\mathbb{T}_n$ and $\slashed{\mathbf{1}}_n$). For algebra $\slashed{\mathbb{C}}$, the set of members of the algebra's uncurling metric group is the same as the set of members of its left regular representation algebra. \item While $\slashed{\mathbb{C}}$ and $\mathbb{R}\oplus\mathbb{R}$, are isomorphic, an isomorphism between any two of $\mathbb{C}$, $\slashed{\mathbb{C}}$, and $\mathbb{D}$ is excluded by Corollary \ref{8/18/23.2} (as indicated in Section 2.3).  
\item  There is no morphism pointing from $\mathfrak{u}_{\mathbb{C}}$ to $\mathfrak{u}_{\mathbb{R}}$. That is, there is clearly no nonzero real matrix $V$ such that $V^T[\sigma]V \subset \left[\begin{array}{ccc}
 \alpha & \beta   \\ 
 \beta & -\alpha    
\end{array}\right]$, since all the real matrix members of $V^T[\sigma]V$  must have determinant zero - while $\left[\begin{array}{ccc}
 \alpha & \beta   \\ 
 \beta & -\alpha    
\end{array}\right]$ has determinant zero if and only if $\alpha,\beta = 0$ and is thereby the zero matrix. This is distinct from the usual demonstration that $\mathbb{C}$ is simple.  On the other hand, there is clearly a morphism from $\mathfrak{u}_{\mathbb{R}\oplus\mathbb{R}}$ to $\mathfrak{u}_{\mathbb{R}}$.
\item There is a morphism between $A_{10}$ and $\mathbb{R}\oplus\mathbb{R}$ associated with the ideal $\left\{\left[\begin{array}{ccc}
 0 & z   \\ 
 0 & 0    
\end{array}\right]:z\in \mathbb{R}\right\}$.  The functor implies that there then must exist a morphism from $\mathfrak{u}_{A_{10}}$ and $\mathfrak{u}_{\mathbb{R}\oplus\mathbb{R}}$.  Indeed, existence of the latter morphism is obvious from the last entries in rows 10 and 4 of Table I. 
\item There is an epimorphism from $A_{13}$ to $\mathbb{R}\oplus\mathbb{R}$ associated with the ideal

 $\left\{\left[\begin{array}{ccc}
0  & z & w   \\
 0 & 0 & v \\
 0 & 0 & 0    
\end{array}\right]:z,w,v\in\mathbb{R}\right\}$. The functor then implies that there must exist a morphism from $\mathfrak{u}_{A_{13}}$ and $\mathfrak{u}_{\mathbb{R}\oplus\mathbb{R}}$.  Indeed, existence of the latter morphism is obvious from the last entries in rows 13 and 4 of Table I.  

On the other hand, as regards $A_{13}$ and its subalgebra $A_{12}$, Corollary \ref{8/18/23.2} indicates that there can be no morphism from $\mathfrak{u}_{A_{13}}$ to $\mathfrak{u}_{A_{12}}$.  This precludes an epimorphism between the respective algebras - a demonstration that does not employ an ideal-based argument.
\item 
There is an epimorphism from $A_{12}$ to  $\mathbb{D}$ associated with the ideal

 $\left\{\left[\begin{array}{ccc}
0  & 0 & w   \\
 0 & 0 & v \\
 0 & 0 & 0    
\end{array}\right]:w,v\in\mathbb{R}\right\}$.
The functor then implies that there must exist a morphism from $\mathfrak{u}_{A_{12}}$ to $\mathfrak{u}_{\mathbb{D}}$.  Indeed, existence of the latter morphism is obvious from the last entries in rows 12 and 5 of Table I. 
\item There is an epimorphism from $A_{11}$ to $\mathbb{D}$ associated with the ideal 

 $ \left\{\left[\begin{array}{ccc}
0  & 0 & w   \\
 0 & 0 & 0 \\
 0 & 0 & 0    
\end{array}\right]:w\in\mathbb{R} \right\} $. Accordingly, there must exist a morphism between $\mathfrak{u}_{A_{11}}$ and $\mathfrak{u}_{\mathbb{D}}$  - and it clearly does by inspection of the depicted uncurling metric groups at the end of rows 11 and 5.  
\end{itemize}



\hspace{-3.5cm} $ 
\begin{array}{cccccccc}
{\rm \underline{algebra}} & \mbox{\underline{usual norm}} & \mbox{\underline{special unital norm family}} & \mbox{\underline{uncurling metric group}} &  \bigskip \\ 
M_n(\mathbb{R}) & ({\rm det}(s))^n & ({\rm det}(s))^{\frac{1}{n}} & \alpha \,(n^2\times n^2)\mbox{\rm -``transpose\, matrix"}  \\  \\ 
\mathbb{C} & x^2+y^2 = r^2 & \sqrt{x^2+y^2}\,e^{\beta{\rm arctan}\left(\frac{y}{x}\right)} = re^{\beta\theta} & \left[\begin{array}{ccc}
 \alpha & \beta   \\ 
 \beta & -\alpha    
\end{array}\right]  =  \slashed{\mathbf{1}}_2 \, \mathbb{C}_\text{left} \\ \\ \slashed{\mathbb{C}}
 & x^2-y^2 & \begin{array}{cc}
 \sqrt{x^2-y^2}\, e^{\beta{\rm arctanh}\left(\frac{y}{x}\right)}   \\
 = (x^2-y^2)^{\frac{1}{2}}\left( \frac{x+y}{x-y}\right)^\frac{\beta}{2}    
\end{array} & \left[\begin{array}{ccc}
 \alpha & \beta   \\
 \beta & \alpha    
\end{array}\right]  = \slashed{\mathbb{C}}_\text{left} \\ \\
\mathbb{R}\oplus\mathbb{R} & xy & (xy)^{\frac{1}{2}}\left(\frac{x}{y}\right)^{\frac{\sigma}{2}} & \left[\begin{array}{ccc}
\sigma_1 & 0   \\
 0 & \sigma_2    
\end{array}\right] = [\mathbb{R}\oplus\mathbb{R}]_\text{left}  \\  \\
\vspace{-1mm} \mathbb{D}  & x^2 & x\, e^{\beta\frac{y}{x}} & \left[\begin{array}{ccc}
 \alpha & \beta   \\
 \beta & 0    
\end{array}\right] = \slashed{\mathbf{1}}_2\mathbb{D}_\text{left} \\ \\

\vspace{-1mm} \mathbb{H}  & (x^2+y^2+z^2+w^2)^2 & \sqrt{x^2+y^2+z^2+w^2} &  \,{\rm diag}\{\alpha,-\alpha,-\alpha,-\alpha\}   \\ \\

\mbox{\rm Cayley-Dixon seq., } \mbox{after }\mathbb{H} & (\mbox{\rm Euclidean norm})^2 & {\rm Euclidean\, Norm} &  \,{\rm diag}\{\alpha,-\alpha,-\alpha,\dots,-\alpha\} &  \\ \\
\mbox{Spin Factor Jordan} & (\mbox{\rm Minkowski norm})^2 & \mbox{\rm Minkowski norm} &  \,{\rm diag}\{\alpha,\alpha,\dots,\alpha\} & \\ \\ \bigoplus_{i=1}^n \mathbb{R}
& \prod x_i & \left(\prod x_i^{\sigma_i}\right)^\frac{1}{n}, \sum\sigma_i=n & {\rm diag} \{\sigma_1,\dots,\sigma_n\} & \\ \\
\left\{ \left[\begin{array}{ccc}
 x & z   \\
 0 & y    
\end{array}\right]: (x,y,z)\in \mathbb{R}^3 \right\} & x^2y & (xy)^{\frac{1}{2}}\left(\frac{x}{y}\right)^{\frac{\sigma}{2}} & \left[\begin{array}{ccc}
\sigma_1  & 0 & 0   \\
 0 & \sigma_2 & 0 \\
 0 & 0 & 0    
\end{array}\right] & \\ \\
\left\{  \left[\begin{array}{ccc}
x  & z & w   \\
 0 & x & z \\
 0 & 0 & x    
\end{array}\right]: (x,z,w)\in \mathbb{R}^3   \right\}& x^3 & x\, e^{ \beta\frac{z}{x}  +\gamma \frac{w}{x} -\gamma \frac{ 1}{2}\left(\frac{z}{x}\right)^2 } &  \left[\begin{array}{ccc}
\alpha  & \beta & \gamma   \\
 \beta & \gamma & 0 \\
 \gamma & 0 & 0    
\end{array}\right]  \\ \\ \left\{ \left[
\begin{array}{ccccc}
x & z & w   \\
0 & x & v \\
0 & 0 & x  
\end{array}
\right] :  (x,v,z,w)\in\mathbb{R}^4 \right\} & x^4 & x\, e^{\beta\frac{v}{x} + \gamma\frac{z}{x} +\delta\frac{w}{x}- \frac{\delta}{2}\left(\frac{zv}{x^2}\right)}
& \left[
\begin{array}{cccc}
\alpha & \beta & \gamma & \delta   \\
\beta & 0 & \frac{1}{2}\delta & 0 \\
\gamma & \frac{1}{2}\delta &  0 & 0   \\
\delta & 0 &  0 & 0
\end{array}
\right]\\ \\
\left\{  \left[\begin{array}{ccc}
x  & z & w   \\
 0 & x & v \\
 0 & 0 & y    
\end{array}\right]:  (x,y,v,z,w)\in\mathbb{R}^5 \right\}  & x^4y & (xy)^{\frac{1}{2}}\left(\frac{x}{y}\right)^{\frac{\sigma}{2}}  e^{\frac{\beta}{2}  \frac{z}{x}}  & \left[\begin{array}{ccccc}
\sigma_1  & 0 & 0 & \beta & 0   \\
 0  & \sigma_2 & 0 & 0 & 0  \\
0   & 0 & 0 & 0 & 0 \\
 \beta  & 0 & 0 & 0 & 0  \\
 0  & 0 & 0 & 0 & 0    
\end{array}\right]  \\  
\mathbb{T}_n & x_1^n & x_1 e^{P_{\gamma_2,\dots,\gamma_n}\left(\frac{x_2}{x_1},\dots,\frac{x_n}{x_1}\right)} & \slashed{\mathbf{1}}_n[\mathbb{T}_n]_\text{left} & \bigskip \end{array} 
$

\bigskip\bigskip


\centerline{\bf Table I}

\newpage

\begin{itemize}\item A simple calculation shows that $\mathfrak{u}_{\mathbb{H}}$ is given by the parametrized matrix $\text{diag}\{\alpha,-\alpha,-\alpha,-\alpha\}$.  This indicates that  $\mathbb{H}$ is simple (Corollary \ref{8/17/22.1}).  This is an alternative to the ideal-based argument that $\mathbb{H}$ is simple because it can be shown to be a division algebra.


 Independent of the above, the functor indicates that there cannot be an epimorphism from  $\mathbb{H}$ to its subalgebra $\mathbb{C}$ since 	their respective uncurling metric groups depicted in row 6 and row 2 preclude a morphism pointing from $\mathfrak{u}_{\mathbb{H}}$ to $\mathfrak{u}_{\mathbb{C}}$, due to Corollary \ref{8/18/23.2}.  This is again an argument that does not invoke ideals.\end{itemize}


The above observations related to Table I all refer to the uncurling metric component, $\mathfrak{u}$, of the inversion characterization group.  However, the unital norm component, $U$, can also play an important role.  For example, a glance at the uncurling metrics for $\slashed{\mathbb{C}}$ versus those of $\mathbb{R}\oplus\mathbb{R}$ fails to immediately indicate that these algebras are intimately related.   In contrast, a glance at the sets of special unital norms for these algebras immediately suggests an intimate relationship between the two algebras. Conversely, a glance at the groups of unital norms (their subset of special unital norm families are shown in Table I) for $\mathbb{C}$, $\slashed{\mathbb{C}}$, and $\mathbb{D}$ immediately
\enlargethispage{1cm} indicates that they cannot be isomorphic (Theorem \ref{6/6/23.3}).  This kind of observation can be expected to be much more useful in higher dimensions.

\subsection{Further classification considerations} For each of the algebras in Table I (but using $n=2$ in row 1 and row 14), uncurling metric groups of type $j\ne 0,1$ are easily computed by hand when  $|j|$ is ``small".  
Thereby, one quickly discovers that within each algebra those groups of different type are identical.  This surprising behavior is of interest in itself.  That is, let $\hat{s}$ be the left regular representation of $s$, so that we can write $s^j = \hat{s}^{j-1}[s]$, where $s$ ranges over the algebra elements. With respect to the basis defined by the columns of an uncurling metric $L$, the above behavior implies that the various $s^j$ are simultaneously irrotational, at least for small $|j|$.  Now consider the vector field $F(s)[s]$ where $F$ is an \mbox{$(n\times n)$-matrix} whose entries are smooth functions of $s$, and suppose $F(s)[s]$ is irrotational.  In general, this does not imply that $(F(s))^j[s]$ is irrotational for $j\ne 0,1$ (e.g., for \mbox{$n=2$}, $s=(x,y)$, and $j=2$, let $F= \left[\begin{array}{ccc}
 f(x) & 1   \\ 
 1 & g(y)    
\end{array}\right]$ with nonconstant $f(x)$, $g(y)$).  So there is something special about the map defined by $\hat{s}$ that results in the unexpected behavior cited above. 
In practical terms, that behavior is of additional interest because it impacts the coarseness of an alternative algebra classification system resulting from our approach (which is dependent on the number of distinct isomorphism invariants that are available, which are presumably dependent on the number of distinct uncurling metric groups).  This important consideration is identified as a matter for future investigation. At this point we are simply left with,
\begin{conjecture} For a given algebra, the uncurling metric groups of all types $j\ne 0,1$ are identical. 
\end{conjecture}
\noindent If the conjecture is false, our alternative algebra classification program is strengthened accordingly.

 From Theorem \ref{6/6/23.3}, each algebra is associated with an invariant derived from the logarithm of the parametrized form of the unital norm group, given by an ordered triple consisting of the number of parameters multiplying the logarithm of a rational function, the number of parameters multiplying the arctangent of a rational function, and the
 \enlargethispage{1cm} number of parameters multiplying a rational function.  For the direct sum of algebras, these triples add component-wise (as is evident when the direct sum is expressed using block matrices).

On the other hand, examples of special unital norm families are listed in \mbox{column 3} of Table I.  These necessarily have one less independent parameter than the unital norm group.  The missing parameter can be thought of as originally associated with the one-dimensional subgroup of $\mathfrak{u}$ consisting of multiples of the uncurling metric that leads to the usual norm or its surrogate (according to Theorem \ref{5/13/22.1} or Corollary \ref{11/17/22.1}).  Since all unital norm groups $U$ have a one-dimensional subgroup associated with such a parameter, in drawing distinctions between algebras it makes sense to suppress it - i.e., instead deal with the special unital norm family.  Specifically, if we consider the subset of $\log U$ composed of logarithms of only the special unital norms, we can classify algebras according to how many independent parameters in that ``log special unital norm family" multiply the logarithm of a rational function, how many independent parameters multiply the arctangent of a rational function, and how many independent parameters multiply a rational function - thereby associating each algebra with an invariant ordered triple of integers differing from the original ordered triple by having 1 subtracted from the value of its first component.
 

For example, from Table I it is thereby seen that $M_n(\mathbb{R})$ is associated with $(0,0,0)$,  $\slashed{\mathbb{C}}$ is associated with $(1,0,0)$, $\mathbb{C}$ is associated with $(0,1,0)$, and $\mathbb{D}$ is associated with $(0,0,1)$.  Such triples are also immediately evident for the other algebras in \mbox{Table I}.
We may consider that the three two-dimensional algebras exemplify in pure form the ``logarithmic", ``angular" (i.e., resulting from an arctangent), and ``rational" class of special unital norm families.  That is, $\mathbb{C}$ is of pure angular type, $\slashed{\mathbb{C}}$ is of pure logarithmic type, and $\mathbb{D}$ is of pure rational type. One might also say that $M_n(\mathbb{R})$ and $\mathbb{H}$ are of null type. 
 
 \section{The case of nonunital algebras} Nonunital algebras can be addressed by modifying Definition \ref{7/21/23.1} to remove the reference to an algebra's units and multiplicative identity, and adding the restriction $j\ge 0$.   However, unital norms are no longer germane.  That is, integration of $Ls^j$, $j\ge 0$, does not have the interest that integration of $Ls^{-1}$ has in the case of unital algebras (it does not lead to anything comparable to Theorem \ref{6/6/23.3}).

For nilpotent algebras, the uncurling metric groups of different types will be different, something that is not evident for the case of unital algebras. For example, consider the real three-dimensional nonunital algebra composed of nilpotent elements $(x,y,z)$ that multiply according to the representation
  $ \left[
\begin{array}{ccccc}
0 & x & y & z  \\
0 & 0 & x & y \\
0 & 0 & 0 & x \\
0 & 0 & 0 & 0
\end{array}
\right] $.  We then have $(x,y,z)^2 = (0,x^2, 2xy)$.  The uncurling metric group resulting from replacement of $s^{-1}$ with  $s^2$ in Definition \ref{7/21/23.1} is then $\left[
\begin{array}{ccccc}
\alpha & \beta & \gamma   \\
\beta & \gamma & 0 \\
\gamma & 0 & 0  
\end{array}
\right] $.  On the other hand, substitution of $s^{-1}$ with $s^3$  deals with $(x,y,z)^3 = (0,0, x^3)$, and the associated uncurling metric group becomes $\left[
\begin{array}{ccccc}
\alpha & \beta & \gamma   \\
\beta & \delta & 0 \\
\gamma & 0 & 0  
\end{array}
\right] $.  Substitution of $s^{-1}$ by $s^j$, $j>3$, results in the additive group of all symmetric $(3\times 3)$-matrices since $s^j$ is zero for all algebra elements.  The uncurling metric groups associated with $s^2$, $s^3$, and \enlargethispage{1cm} $s^4$, must differ because the respective dimensions of the ranges of these operations differ - and this type of behavior will be typical with nilpotent algebras.  
 

\bigskip

\appendix

 \section{First motivation: A  new proof of the Pythagorean Theorem}    The isotropy of the Newtonian concept of space, asserting that physical laws are independent of direction (which results in ``conservation of angular momentum" as an example of the renowned ``Noether's Theorem" in theoretical physics), has a non-physical analogue with respect to Euclidean space or, for our present purposes, the Euclidean plane.  Specifically, given a line and a point, the ``Parallel Postulate" (or its equivalent) ultimately indicates that not only is there is a unique line through the given point parallel to the given line, but also that there is a unique line though the given point perpendicular to the given line.  Using the formulation of Euclidean geometry due to George Birkhoff \cite{birkhoff:1932}, which fundamentally incorporates $\mathbb{R}$, it is then easy to derive a bijection between the points of the Euclidean plane and those of $\mathbb{R}\times\mathbb{R}$, using a pair of perpendicular lines intersecting at some chosen point $O$, such that this ``origin" $O$ is mapped to $(0,0)$.  That is, the perpendicular lines through $O$ then become directed ``orthogonal coordinate axes": the real-valued coordinates are due to the assertion of one of Birkhoff's Euclidean geometry postulates that there is a bijection between the points of $\mathbb{R}$ and the points of any line, which also endows the line with a direction as determined by the increasing values of $\mathbb{R}$.  Consequently, this formulation allows incorporation of concepts of real vector spaces, limits, continuity, and one-dimensional Differential Calculus.  Thereby, the above bijection (which was specifically demonstrated in \cite{birkhoff:1932}) leads to an isometric isomorphism. 
To see this, we specify that $\mathbb{R}\times\mathbb{R}$  inherits the Euclidean length function $\ell(P)$, where the latter gives the length of the line segment with endpoints $O$ and $P$, where $P$ is  any point in the Euclidean plane.  Accordingly, we specify that $\ell(x,y) \equiv \ell(P)$ for $(x,y)$ the image of $P$ under the bijection.  Furthermore, the Euclidean plane with distinguished origin $O$ is easily made into a real vector space (dilation/reflection of directed line segments based at $O$ result from multiplication by members of $\mathbb{R}$, and the sum of directed line segments based at $O$ follow from parallelogram constructions).   Since $\mathbb{R}\oplus\mathbb{R}$ is already a real vector space, and both it and the latter Euclidean plane vector space are endowed with the above length function (a degree-1 positive homogeneous function, as implied by one of the geometry postulates), the isometric isomorphism follows without difficulty (this being simply a formalization of Descartes' conception).  An important feature is that (due to the Parallel Postulate or equivalent) such an isomorphism can be constructed from {\it any} pair of perpendicular lines through $O$, i.e., regardless of the angle one of the perpendiculars makes with an arbitrarily selected line through $O$.  So this is the isotropic feature of the Euclidean plane referred to at the beginning of this appendix.    
 
 We will now present a new proof of the Pythagorean Theorem based on that isotropy (fully detailed in \cite{fgreensite:2022b}), and it will be seen that features of the derivation motivate Definition \ref{7/21/23.1}.
 
 First,  we have the usual directional derivative of the inherited length function on $\mathbb{R}\oplus\mathbb{R}$, denoted $D\ell(x,y;u,v)$, this being the value of what is essentially a one-dimensional derivative of the length function at $(x,y)$ with respect to the direction vector $(u,v)\ne (0,0)$ (i.e., employing the limit of an expression involving the usual vector space operations on a one-dimensional subset of $\mathbb{R}\oplus\mathbb{R}$).  For $(x,y)\ne (0,0)$, this derivative value can be shown to exist based on the triangle inequality. The analogous manipulations in the above Euclidean plane vector space allow computation of a derivative similarly notated as $D\ell(P;C)$.  If $P$ and $C$ are the respective pre-images of $(x,y)$ and $(u,v)$ under the isomorphism resulting from a set of orthogonal coordinate axes in the Euclidean plane vector space, then $D\ell(P;C) = D\ell(x,y;u,v)$ (because the isomorphism is isometric).  With respect to a given direction vector, it can be shown that $\ell(x,y)$ is continuously differentiable except at the origin (this being the only subtlety of the proof, since the two dimensional continuity of the derivative must be established in the absence of an expression for the length of an arbitrary line segment; the continuity demonstration ultimately results from the degree-1 positive homogeneity of the length function, the latter essentially being an expression of one of the Euclidean geometry postulates).  Using the Mean Value Theorem of one-dimensional Calculus, it is then easy to show \begin{equation}\label{11/30/23.1} D\ell(x,y;u,v) = \frac{\partial \ell(x,y)}{\partial x}u + \frac{\partial \ell(x,y)}{\partial y}v,\end{equation} where $\frac{\partial \ell(x,y)}{\partial x}$ and $\frac{\partial \ell(x,y)}{\partial y}$ are the directional derivatives of $\ell$ at $(x,y)\ne(0,0)$, specified by the direction vectors $(1,0)$ and $(0,1)$ on the respective coordinate axes in $\mathbb{R}\oplus\mathbb{R}$, where the latter direction vectors are the images of unit length direction vectors associated with a chosen pair of orthogonal coordinate axes in the Euclidean plane vector space that result in a first isometric isomorphism as indicated above.  The latter coordinate axes are respectively labeled as the $x$-axis and the $y$-axis. We conclude from (\ref{11/30/23.1}) that, for fixed $(x,y)\ne (0,0)$,  $D\ell(x,y;u,v)$ is a linear functional, and is evidently the dual of $\left(\frac{\partial\ell(x,y)}{\partial x}, \frac{\partial\ell(x,y)}{\partial y}\right)$.  From that first isomorphism, the vector expression at the end of the prior sentence labels a point $G$ of the Euclidean plane vector space.  An alternative orthogonal coordinate axis system through the origin in the Euclidean plane, resulting in a second isometric isomorphism with $\mathbb{R}\oplus\mathbb{R}$, will provide a different label for this {\it same} point $G$ following computation of a linear functional analogous to (\ref{11/30/23.1}) expressed via that second isomorphism, because each orthogonal coordinate system provides an isometric isomorphism with $\mathbb{R}\oplus\mathbb{R}$ thereby allowing computation of the directional derivative $D\ell(P;C)$ for the Euclidean plane vector space.  In fact, we can find a set of alternative orthogonal coordinate axes, denoted as the $x'$-axis and $y'$-axis, such that $P$ is on the positive portion of the $x'$-axis - so that the associated isomorphism now labels $P$ as $(x',0)$ instead of $(x,y)$.  The Euclidean vector space point $G$ is now labeled by  
 \begin{equation}\label{11/24/23.4} \left( \frac{\partial \ell(x',0)}{\partial x'}, \frac{\partial \ell(x',0)}{\partial y'}\right) = (1,0),\end{equation} where the right-hand-side is an easy calculation since determination of the two directional derivatives on the left-hand-side only involves evaluation of $\ell$ on the coordinate axes.  Equation (\ref{11/24/23.4}) states that $ \left( \frac{\partial \ell(x',0)}{\partial x'}, \frac{\partial \ell(x',0)}{\partial y'}\right)$ is the point of the unit circle on the ray from $(0,0)$ that includes the point $(x',0)$.
 Based on this second isomorphism, this means that $G$ is the point of the unit circle on the ray from $O$ that includes the point $P$.  Hence, based on the first isomorphism, this requires $\left( \frac{\partial \ell(x,y)}{\partial x}, \frac{\partial \ell(x,y)}{\partial y}\right)$ to be the point of the unit circle on the ray from $(0,0)$ that goes through the point $(x,y)$. Consequently, \begin{equation}\label{11/25/23.1} \ell\left( \frac{\partial \ell(x,y)}{\partial x}, \frac{\partial \ell(x,y)}{\partial y}\right) = 1.\end{equation} and, taking into account that $\ell$ is a degree-1 positive homogeneous function,  \begin{equation}\label{11/24/23.5} \left( \frac{\partial \ell(x,y)}{\partial x}, \frac{\partial \ell(x,y)}{\partial y}\right) = \left( \frac{x}{\ell(x,y)},  \frac{y}{\ell(x,y)}\right),\end{equation} since the right-hand-side of (\ref{11/24/23.5}) is (uniquely) both a point on the unit circle while also being on the ray from $(0,0)$ that includes $(x,y)$ - as are now seen to be required for a labeling of $G$ via an isometric isomorphism. 
 
  As regards (\ref{11/24/23.5}), one-dimensional indefinite integration (anti-differentiation) can be applied to the equality of the first coordinates, and then one-dimensional indefinite integration can be independently applied to the equality of the second coordinates.  Reconciliation of the  two resulting expressions for $\ell(x,y)$ in the context of (\ref{11/25/23.1}) and the nonnegativity of $\ell$ quickly leads to $\ell(x,y)=\sqrt{x^2+y^2}$, from which the Pythagorean Theorem easily follows via an argument involving parallelograms to measure lengths of arbitrary line segments using $\ell$. 
   The full details of this proof are found in \cite{fgreensite:2022b}.

  
So now we begin a series of easy generalizations.  

Equations (\ref{11/24/23.5}) and (\ref{11/25/23.1}) can be recast as, \begin{equation}\label{11/29/23.5} s = \ell(s)\nabla\ell(s),\end{equation} and\begin{equation}\label{11/29/23.6}\ell\left(\nabla\ell(s)\right) = 1,\end{equation} now generalized to $s\in \mathbb{R}^n$, $n\ge 2$, where $\nabla\ell(s)$ is simply notation for the $n$-tuple of coordinate-wise derivatives with respect to the standard basis, as on the left-hand-sides of  (\ref{11/25/23.1}) and (\ref{11/24/23.5}).  For nonnegative $\ell(s)$, the solution to (\ref{11/29/23.5}) and (\ref{11/29/23.6}), is similarly determined to be the Euclidean norm.  

But why stop there?  Equations (\ref{11/29/23.5}) and (\ref{11/29/23.6}) are amenable to further generalization, from which  we can obtain {\it any} quadratic space, via the norm $\ell(s)$ \mbox{specified by} \begin{equation}\label{11/29/23.1} Ls = \ell(s)\nabla\ell(s),\end{equation} and \begin{equation}\label{11/29/23.2}\ell\left(L^{-1}\nabla\ell(s)\right) = 1,\end{equation} for any appropriately-dimensioned nonsingular symmetric $L$.  The matrix $L$ must be symmetric since the right-hand-side of (\ref{11/29/23.1}) is the gradient of $\frac{1}{2}\ell^2(s)$, so $Ls$ is required to be a gradient.  Equation (\ref{11/29/23.2}) follows from dividing both sides of (\ref{11/29/23.1}) by $\ell(s)$, then applying $L^{-1}$ to both sides of the resulting equation, then applying $\ell$ to both sides of that resulting equation, and finally invoking the assumed degree-1 positive homogeneity of $\ell(s)$. Again,  $\nabla\ell(s)$ is simply notation for the ordered $n$-tuple of  coordinate-wise one-dimensional derivatives, with the standard basis of $\mathbb{R}^n$ assumed.
  
Before getting into generalization of the above to algebras, let us examine the things a bit further.  Equations (\ref{11/29/23.1}) and (\ref{11/29/23.2}) are feasible because $Ls$ is a gradient and the equations continue to hold when $s$ is replaced by $\alpha s$, $\alpha >0$.  This is a consequence of the Euler Homogeneous Function Theorem, which states that if $\ell(s)$ is a degree-1 positive homogeneous function then
  $s\cdot\nabla \ell(s) = \ell(s)$.  For $\alpha>0$, replacing $s$ by $\alpha s$ then leads to \[ \alpha s \cdot \nabla \ell(\alpha s) = \ell(\alpha s) = \alpha \ell(s) = \alpha s\cdot \nabla \ell(s),\] where $\nabla \ell(\alpha s)$ means that $\nabla \ell$ is evaluated at $\alpha s$. Equality of the above left-hand-side expression with the above right-hand-side expression implies, \begin{equation}\label{11/30/23.5} \nabla \ell(\alpha s) = \nabla \ell(s).\end{equation}    
  Given that we require a degree-1 positive homogeneous $\ell(s)$, this means that the left-hand-side of (\ref{11/29/23.1}) and the right-hand-side of (\ref{11/29/23.1}) are consistent when $s$ is replaced by $\alpha s$, $\alpha>0$ (each side simply being multiplied by $\alpha$).  
  
  With that in mind, suppose we are now given a real unital associative algebra with $\mathbb{R}^n$ as its vector space of elements  and, inspired by the above norms on vector spaces, we wish to look for a degree-1 positive homogeneous function $\ell(s)$ that somehow also acknowledges the algebra.  The degree-1 positive homogeneity will again require satisfaction of $\nabla\ell(\alpha s) = \nabla\ell(s)$.  Along those lines, we recall that the ``vector-space-appropriate" (\ref{11/29/23.1}) exhibits the feature that both sides of the equation are multiplied by a common factor when $s$ is replaced by a multiple of $s$.  So, we want an ``algebra-appropriate" equation that looks similar to (\ref{11/29/23.1}), and retains the feature noted in the prior sentence.  An obvious candidate is simply, \begin{equation}\label{12/22/23.1} Ls^{-1}=\frac{1}{\ell(s)}\nabla\ell(s),\end{equation} since the above looks a lot like (\ref{11/29/23.1}), it addresses the algebra through its use of the nonlinear unary operation $s^{-1}$, and replacement of $s$ by $\alpha s$ leads to both sides of the equation being multiplied by a common factor (this time it is \mbox{$\frac{1}{\alpha}$, instead of $\alpha$)}.   
  
  But how do we augment (\ref{12/22/23.1}) to assure degree-1 positive homogeneity of $\ell(s)$?  Suppose we require $\ell(s)$ to also respect the algebra's multiplicative inversion operation, \begin{equation}\tag{\ref{5/23/23.3}} \ell(s^{-1}) = (\ell(s))^{-1} = \frac{1}{\ell(s)}.\end{equation}   In that case, (\ref{12/22/23.1}) becomes \begin{equation}\label{11/27/23.7} Ls^{-1} = \ell(s^{-1})\nabla\ell(s).\end{equation}  Applying $L^{-1}$ to both sides of (\ref{11/27/23.7}), then applying $\ell$ to both sides of the resulting equation, and then utilizing the degree-1 positive homogeneity of $\ell(s)$, we obtain
\begin{equation}\tag{\ref{11/29/23.2}}\ell\left(L^{-1}\nabla\ell(s)\right) = 1.\end{equation}  Conversely, it is also clear that (\ref{11/29/23.2}) and (\ref{11/27/23.7}) cannot both hold unless $\ell(s)$ is degree-1 positive homogeneous.
  
  Thus, we are motivated to find $\ell(s)$ satisfying (\ref{11/27/23.7}) and (\ref{11/29/23.2}).
  
  However, there is a new constraint necessarily introduced in this extension to unital algebras.  We began with (\ref{12/22/23.1}), whose right-hand-side is $\nabla\log\ell(s)$, i.e., a gradient.  In general, $s^{-1}$ is not a gradient, nor is $Ls^{-1}$ with an arbitrary symmetric $L$. Consequently, to proceed as above, it is now required that $L$ satisfies \[ d([Ls^{-1}]\cdot ds) = 0.\]  It might be speculated that, for any unital algebra with vector space space of elements $\mathbb{R}^n$, the set of $L$ satisfying this constraint always has nontrivial members, and the set often distinguishes different algebras from each other.  It might also be speculated that, given a member $L$ of this set,  the coupled equations, (\ref{11/29/23.2}) and (\ref{11/27/23.7}), emulating (\ref{11/25/23.1}) and (\ref{11/24/23.5}), always has a solution if $L$ is nonsingular.
  
  \enlargethispage{1cm}
  
These speculations can be taken to be the origin of Definition \ref{7/21/23.1} with $j=-1$, the beginning of our program.
  
  \bigskip
  
  \section{Second motivation: an unusual regularization method}
     Appendix A inverts the reality that, rather than motivating our program, the forwarded new proof of the Pythagorean Theorem was instead itself motivated by our program (see Remark \ref{11/21/23.1}). In this Appendix B, we will present what actually led to the introduction of uncurling metrics and unital norms.  Instead of relying on Euclidean geometry and some Calculus as in Appendix A, this motivation relies on the interplay of deterministic and statistical interpretations of an unusual regularization scheme in the realm of Inverse Problems (essentially, a geometrical interpretation of a statistically-inspired method) \cite{fgreensite:2022d}.
  
    
 Consider the linear problem described by, \begin{equation}\label{7/7/20.9}  y^\delta=Fx + \nu,\end{equation} where one is given a data vector $y^\delta$ that is a version of an unknown source vector $x$ as transformed by the known linear operator $F:\mathbb{R}^m\rightarrow\mathbb{R}^n$, but where the data vector has been corrupted by a ``noise" vector $\nu$ that is only characterized by \begin{equation}\label{7/12/20.1} \| y^\delta - Fx\|^2 \le \delta^2,\end{equation} where $\delta$ is given and $\|\cdot \|$ is the Euclidean norm.  
    
   If $F$ is an ill-conditioned matrix (e.g., a discretization of a linear compact operator), a solution estimate $F^{-1}y^\delta$ or $F^\dagger y^\delta$ is usually worthless ($F^\dagger$ is the pseudoinverse of $F$).  That is, consider the singular system for $F$ as  $\{u_i ,\sigma_i, v_i\}$, where $\{u_i\}$ is the set of eigenvectors of $FF'$, $\{v_i\}$ is the set of eigenvectors of $F'F$, and $\{\sigma_i^2 \}$ is the set of eigenvalues of $F'F$ (i.e., $\{\sigma_i \}$ is a set of singular values).   The least squares solution estimate of smallest norm is then, \begin{equation}\label{7/7/20.8} x_{\rm poor} = F^\dagger y^\delta = \sum_{ \{i: \sigma_i \ne 0\}} \frac{u_i' y^\delta}{\sigma_i}v_i,\end{equation} where  an expression like $a' b$ indicates the Euclidean inner product of vectors $a$ and $b$.   It is evident that the existence of small singular values  leads to potentially extreme noise amplification, so that the above is a poor estimate of $x$ as regards the at least possibly approachable $x^\dagger \equiv F^\dagger F x$.
 Instead, a common treatment (upon which zero-order Tikhonov regularization with the ``Discrepancy Principle" is based \cite{vogel:2002,engl:2000}) is to supply the solution estimate of minimum norm that is consistent with the modification of (\ref{7/12/20.1}) as \begin{equation}\label{7/7/20.3} \| y^\delta - Fx\|^2 =\delta^2 .\end{equation}   That is, the solution estimate is the point of smallest Euclidean norm on the ``discrepancy ellipsoid",   \begin{equation}\label{7/15/20.1} \partial\Omega_\delta \equiv \{ x\in \mathcal{D}(F): \| y^\delta - Fx\|^2 =\delta^2\},\end{equation} where $\mathcal{D}(F)$ is the domain of $F$.  We use $\mathcal{D}(F)$ here instead of $\mathbb{R}^m$ because we will want (\ref{7/15/20.1}) to also be relevant when $F$ is a compact operator (the general Inverse Problems setting \cite{engl:2000}).
  A Lagrange multiplier argument then gives the solution estimate as, \begin{equation}\label{7/8/20.9} x^\delta_D = (F'F + \gamma I)^{-1}F'y^\delta,\end{equation} where $I$ is the identity operator, and regularization parameter $\gamma$ is the reciprocal of the
   Lagrange multiplier as determined by the requirement that (\ref{7/7/20.3}) be satisfied (the subscript ``$_D$" indicates the estimate follows from a ``deterministic" treatment).  
 
 Alternatively, one may consider both $x$ and $\nu$ to be zero-mean Gaussian random vectors (``signal" and ``noise", respectively), with the covariance matrix of $x$ being $C_x$ (assumed for the moment to be nonsingular), and the covariance matrix of noise vector $\nu$ being $\epsilon^2 I$ ($\delta^2$ in (\ref{7/7/20.3}) might now be viewed as an estimate of the sum of the autocovariances of the components of $\nu$).  The {\it maximum  a posteriori} (Bayesian) estimate  is then, \begin{equation}\label{7/8/20.5} x^\delta_S = (F'F + \epsilon^2 C_x^{-1})^{-1} F'y^\delta,\end{equation}  where the subscript ``$_S$" indicates that the estimate follows from  a ``statistical" treatment \cite{vogel:2002}. 
  
 We can thus associate the deterministic solution estimate (\ref{7/8/20.9}) with a statistical interpretation (\ref{7/8/20.5}), as resulting from assigned signal and noise covariance matrices each proportional to the identity matrix, with $\gamma$ as the inverse of the square of the  ``signal-to-noise" ratio (this inverse square being the ratio of the respective noise and signal component variances).  Thus, the diagonal entries of the noise covariance matrix are assumed to be identical, and these represent the noise variance $\epsilon^2$ in each of the components of the data vector (i.e., the variances of each component of random vector $\nu$).   So the full noise covariance matrix is assumed known {\it a priori}, perhaps based on knowledge of $\delta^2$ under the supposition that $\delta^2 \approx n\epsilon^2$ in (\ref{7/7/20.3}), where $n$ is the dimension of the range of $F$.  The factor multiplying the identity matrix to give $C_x$, so that this factor is the (assumed uniform) variance of each component of the signal random vector $x$, then follows once $\gamma$ is computed.  Hence, from a statistical standpoint,  (\ref{7/7/20.3}), (\ref{7/8/20.9}) can be re-interpreted as implying a method resulting in {\it simultaneous} estimation of $x$ and the presumably uniform signal component variances $\epsilon^2/\gamma$. 
 
But as long as one is simultaneously estimating the solution and a covariance matrix feature, why stop there?   Instead of assuming that the   signal covariance matrix is proportional to the identity and estimating the proportionality constant simultaneously with the signal, one might only assume that the signal covariance matrix is diagonal in the basis defined by the eigenvectors of $F'F$, and seek to estimate the diagonal elements simultaneously with the signal (i.e., one assumes  that the matrix whose columns are the right singular vectors of  $F$ diagonalizes the signal covariance matrix).
The latter basis, $\{v_i\}$, is the natural one to be used in the context of this type of assumption since, considering the underlying setting where $F$ is a compact operator, one thing we {\it know} is that for the inverse problem to make sense the signal must be in the domain of $F$ (the problem is {\it defined} by $F$ operating on the signal). That is, we already know that the sum of the squares of the signal components in this basis is finite - and (absent other prior information) this is the only basis for which we know that to be true.  So, the assumption is simply that the components of random vector $x$ in this basis are independent, rather than additionally assuming that they all have the same variance (which would correspond to the assumption that the signal covariance matrix is proportional to $I$).

Proceeding in this way, one notes that the value of the $i$-th eigenvalue of $C_x$ is then the autocovariance of the $i$-th component of the signal random vector in the $\{v_i\}$-basis, i.e., the expectation of the square of the $i$-th component of random vector  $x$ in that basis (which is the expectation of $(v_i'x)^2$).   Thus, an alternative approach to the statistical treatment of (\ref{7/7/20.9}), (\ref{7/7/20.3}) (which we do not claim to be optimal in the solitary estimation of $x$)  would be to require that the eigenvalues of $C_x$ be the squares of the components of the solution estimate vector -  inspired by the actual meaning of those covariance matrix values as expectations  of the squares of the signal component random variables in the $\{v_i\}$-basis.  Such a solution estimate would then be  a fixed point of the dynamical system on $\mathbb{R}^m$ defined by iteration of the operator, \begin{equation}\label{10/19/10.3} A_\epsilon[w] \equiv \lim_{\alpha\rightarrow 0}\left(F'F+\epsilon^2\, V\mbox{\rm diag}\{(v_i'w)^2 + \alpha^2\}^{-1}V' \right)^{-1} F'y^\delta, \end{equation} where $V$ is the matrix whose $i$-th column is $v_i$ for each $i$. The role of $\alpha$ and $\lim_{\alpha\rightarrow 0}$ is to account for the possibility that one of the $v_i'w$ is zero.  To get the iterations started,  a solution estimate candidate $w$ (even chosen arbitrarily) is associated with a signal covariance matrix candidate $V \mbox{\rm diag}\{(v _i'w)^2\}V'$.    According to (\ref{7/8/20.5}),  $A_\epsilon[w]$ is then an updated solution estimate using the latter signal covariance matrix.  For the next iteration, the squares of the components of this updated solution estimate  $A_\epsilon[w]$ in the $\{v_i\}$-basis are used to define the entries of an updated version of the diagonal matrix on the right-hand-side of (\ref{10/19/10.3}) (i.e., updating the signal covariance matrix), which then supplies a successive updated solution estimate as $A_\epsilon[A_\epsilon[w]]$, etc. This iteration process can be repeated indefinitely.  A fixed point of the dynamical system evidently supplies a covariance matrix whose eigenvalues are the squares of the solution estimate components in the $\{v_i\}$-basis.  
 The fixed point $x=0$ is associated with a signal covariance matrix as the matrix of zeros. Thus, one would not select $w=0$ as an initial guess for the solution estimate to begin the iteration process (also, $w=0$ is very far from satisfying (\ref{7/7/20.3})).  
 
 
But it is helpful that one can easily identify the fixed points of $A_\epsilon$. Inserting the previously noted
singular system associated with $F$ into the right-hand-side of (\ref{10/19/10.3}),  we obtain \begin{equation}\label{3/31/12.3} A_\epsilon[w] = \sum_i \frac{\sigma_i (u_i'y^\delta)}{\sigma_i^2 + \epsilon^2/(v_i'w)^2}v_i .\end{equation} A fixed point $p$ must satisfy  \begin{equation}\label{4/1/12.1} A_\epsilon[p] = \sum_i  \frac{\sigma_i (u_i'y^\delta)}{\sigma_i^2 + \epsilon^2/(v_i'p)^2}v_i = p = \sum_i  (v_i' p)v_i.\end{equation}  Since the coefficients of $v_i$ in the two sums in (\ref{4/1/12.1}) are necessarily equal, for $v_i'p\ne 0$ we obtain \begin{equation}\label{7/11/10.1} \sigma_iv_i'p = \frac{u_i' y^\delta \pm \sqrt{(u_i'y^\delta)^2 -4\epsilon^2}}{2}.\end{equation}  In the first place, this means that the series for a fixed point will only contain terms $i$ such that $(u_i'y^\delta)^2 >4\epsilon^2$.  Second, using Taylor series, it is easy to show that if the sign in front of the radical on the right-hand-side above is not the same as the sign of $u_i'y^\delta$, then $|\sigma_i v_i'p| =  O(\epsilon^2/|u_i'y^\delta|)$.  Indeed, whatever signal
 was contained in $u_i'y^\delta$ will have then been mostly 
 removed from the right-hand-side of (\ref{7/11/10.1}) (at least until $|u_i'y^\delta|$ is on the order of $\epsilon$, at which point in practical settings the data is mostly noise-dominated).   Thus, any fixed point of interest will be such that the sign in front of the radical is the same as the sign of $u_i'y^\delta$ .   The solution estimate that takes advantage of the greatest number of data subspaces available to this method is then,
\begin{equation}\label{7/19/10.1} p = \frac{1}{2}\sum_{\{i:\,|u_i'y^\delta|>2\epsilon\}} \left(u_i' y^\delta + {\rm sgn}(u_i' y^\delta) \sqrt{(u_i'y^\delta)^2 -4\epsilon^2}\right)\frac{v_i}{\sigma_i}.\end{equation}      
Note that the fixed points of of $A_\epsilon$ are truncations - a feature this solution estimate shares with the (quite different)  commonly used Truncated Singular Value Decomposition (TSVD) solution estimate \cite{vogel:2002}.  

The following result answers an obvious question of interest, though it is tangental to our objective in this appendix.

\begin{lemma} \label{7/6/21.9} Consider the solution estimate (\ref{7/19/10.1}) but with removal of the terms of the sum on the right-hand-side for which $|u_i'y^\delta|<4\epsilon$.  This is an attracting fixed point of the dynamical system resulting from repeated iterations of $A_\epsilon$. \end{lemma}

The proof is found in \cite{fgreensite:2022d}.

For the purposes of this appendix, the point in presenting all this is that the solution estimate $p$ given by (\ref{7/19/10.1}) is a critical point of the {\it geometric mean} of the solution estimate component absolute values on an implied ellipsoid  derived from the discrepancy ellipsoid on which $p$ lies - in contrast to the most commonly used solution estimate, which is a critical point of the Euclidean norm, i.e., a critical point of the {\it quadratic mean}, on the discrepancy ellipsoid on which it lies.  That is, the geometric mean has both a statistical and deterministic interpretation in the context of the regularization of inverse problems, just as the Euclidean norm does. 
 
 \begin{theorem} \label{7/6/21.6} Consider \begin{equation} \tag{\ref{7/7/20.9}} y^\delta = F x +\nu, \end{equation} where $y^\delta$ is known, $F:\mathbb{R}^k\rightarrow\mathbb{R}^k$ is a given nonsingular linear transformation associated with singular value decomposition matrix $U\Sigma V'$, and $\{u_i\}$ and $\{v_i\}$ are the sets of columns of $U$ and $V$, respectively.  Further consider the dynamical system on $\mathbb{R}^k$ defined by iteration of the operator, \begin{equation}\tag{\ref{10/19/10.3}} A_\epsilon[w] \equiv \lim_{\alpha\rightarrow 0}\left(F'F+\epsilon^2\, V\mbox{\rm diag}\{(v_i'w)^2 + \alpha^2\}^{-1}V' \right)^{-1} F'y^\delta. \end{equation} 
As a solution estimate for $x$, the fixed point $p_\epsilon$ of largest Euclidean norm of this dynamical system  has the following interpretations:
\begin{enumerate}\item {\rm Statistical:} If $\epsilon^2$ is the variance of zero-mean white Gaussian noise random vector $\nu$, then 
$p_\epsilon$ is the  {\rm maximum a posteriori} estimate for zero-mean Gaussian signal random vector $x$ of greatest Euclidean norm
 that results from a signal covariance matrix  whose eigenvectors are the members of the basis $\{v_i\}$ and with the additional feature that its eigenvalues are the squares of the components of the solution estimate with respect to those basis vectors.  \item  
{\rm Deterministic:} Denote the discrepancy of $p_\epsilon$ as $\delta_\epsilon \equiv \| y^\delta - Fp_\epsilon\|$, which may be adjusted by varying $\epsilon$ so as to allow for the influence of prior knowledge concerning the magnitude of unknown noise vector $\nu$.  Let $\mathbb{V}_\epsilon \subset \mathbb{R}^k$ be the space spanned by the subset of eigenvectors $\{v_i: |u_i'y^\delta| > 2\epsilon \}$.  
 Consider the ellipsoid \[ \partial\Omega \equiv \{ s\in \mathbb{V}_\epsilon: \| y^\delta - Fs\|^2 = \delta_\epsilon^2 \}.\]   Then $p_\epsilon\in \mathbb{V}_\epsilon$ is a critical point of the geometric mean of the absolute values of the coordinates of the points of $\partial\Omega$ with respect to the above eigenvector subset basis.
\end{enumerate}
\end{theorem}

\begin{proof}
The statistical interpretation follows from the prior observations related to the form of the right-hand-side of (\ref{7/8/20.5}) and inspection of the form of (\ref{10/19/10.3}) as regards a fixed point.

We now consider the deterministic interpretation.   
To simplify our notation, without loss we assume an arrangement of the sequence of right singular vectors of $F$ such that $\{v_i\}_{i=1}^n$ is the set $\{v_i: |u_i'y^\delta| > 2\epsilon \}$, and we will write $p_\epsilon$ as simply $p$.  We observe that $p$ is a member of $\mathbb{V}_\epsilon$, since we have previously shown that $p$ is given by (\ref{7/19/10.1}). 
Define $\hat{p} \equiv (v_1'p,\dots,v_n'p)$, i.e., $\hat{p}$ is the coordinate expression of $p$ as a member of $\partial\Omega\subset \mathbb{V}_e$ (thus, $\hat{p}$ is $p$ with suppression of the latter's null-valued coordinates with respect to the full basis of $\mathbb{R}^k$ given by $\{v_i\}$). 
Define $C_p \equiv \mbox{diag}\{(v_i'p)^2\}$ and  $C_{\hat{p}} \equiv  \mbox{diag}\{(v_i'p)^2\}_{i=1}^n$.  Since $p$ is a fixed point of $A_\epsilon$,  (\ref{10/19/10.3}) implies \begin{equation} p = \lim_{\alpha\rightarrow 0} \left(F'F+\epsilon^2\,\left[C_p + \mbox{diag}\{\alpha^2\}\right]^{-1}\right)^{-1} F'y^\delta.\label{12/14/20.1}\end{equation} However, the null coordinates of $p$ correspond to the null diagonal values of $C_p$.  There is no intrinsic loss in suppressing these null entries.  {\it Hence, the right-hand-side of (\ref{12/14/20.1}), can be cast as the solution to the variational problem of finding the point minimizing $s'C^{-1}_{\hat{p}} s$ over $z\in \partial\Omega \subset \mathbb{R}^n$, which is $\hat{p}$}.
Specifically, (\ref{12/14/20.1}) implies, and is replaced by, 
\begin{equation} \hat{p} = \left(\hat{F}'\hat{F}+\epsilon^2\,C_{\hat{p}}^{-1}\right)^{-1} \hat{F}'y^\delta, \label{12/19/20.1}\end{equation} where $\hat{F}$ is the operator resulting from the restriction of the domain of $F$ to $\mathbb{V}_\epsilon$.
  According to the principle of Lagrange multipliers, the point minimizing $s'C^{-1}_{\hat{p}} s$ on \[ \partial\Omega = \{s\in \mathbb{V}_\epsilon: \| \hat{F} s -y^\delta\|^2 = \delta_\epsilon^2\} = \{s\in \mathbb{V}_\epsilon: (s'\hat{F}' - (y^{\delta})')(\hat{F} s - y^\delta) = \delta_\epsilon^2\},\] is the critical point of, \[ \mathcal{L}(s,\lambda) \equiv \lambda\left[ (s'\hat{F}' - (y^{\delta})')(\hat{F} s - y^\delta) -\delta_\epsilon^2\right] + s'C^{-1}_{\hat{p}} s.\] 
Setting to zero all of the partial derivatives of $\mathcal{L}$ with respect to the components of $s$ leads to the solution being given as the right-hand-side of (\ref{12/19/20.1}), where the factor $\epsilon^2$ is the reciprocal of the Lagrange multiplier.  The latter is determined by the final equation $\frac{\partial \mathcal{L}}{\partial \lambda} = 0$, which places the solution on $\partial\Omega$.  This solution is indeed $\hat{p}$ since we are {\it given} (\ref{12/19/20.1}) (it is implied by (\ref{12/14/20.1})) and we know that $\hat{p}$ lies on $\partial\Omega$. 

Now define the function  $h:\mathbb{V}_\epsilon\rightarrow \mathbb{R}$ with $h(s) \equiv s'C^{-1}_{\hat{p}} s$. According to the prior paragraph, $\hat{p}$ minimizes $h(z)$ on  $\partial\Omega$.   Consider the level set of $h$ given by the ellipsoid, \[ \mathcal{H} \equiv \{s\in \mathbb{V}_\epsilon: s'C_{\hat{p}}^{-1} s = n \} = \left\{s\in\mathbb{V}_\epsilon: \sum_{i=1}^n \frac{s_i^2}{p_i^2} = n\right\},\] where $s_j\equiv v_j's$ and  $p_j\equiv v_j'p$ for $v_j\in \{v_i\}_{i=1}^n$.  
The point $\hat{p}$ is evidently a member of $\mathcal{H}$.  The tangent space of ellipsoid $\mathcal{H}$ at $\hat{p}$  has normal vectors proportional to the gradient of $h$ at $\hat{p}$, and with respect to the $\{v_i\}_{i=1}^n$-basis,
\begin{equation}\label{3/18/12.1}  \nabla h(s)\Big|_{s=\hat{p}} = \,\, 2\left(\frac{1}{p_1},\dots, \frac{1}{p_n}\right).\end{equation} But as we have noted above, $\hat{p}$ is the solution estimate minimizing $s'C^{-1}_{\hat{p}} s$ on $\partial\Omega$.  Thus, $\mathcal{H}$ is tangent to $\partial\Omega$ at $\hat{p}$.  
 
We next introduce  $g:\mathbb{V}_\epsilon\rightarrow \mathbb{R}$ with $g(s) = \prod_{i=1}^n |s_i|$ (the $n$-th power of the geometric mean of the absolute values of a point's components with respect to the $\{v_i\}_{i=1}^n$-basis).  The level set of $g$ that contains $\hat{p}$ is,  \[\mathcal{G} \equiv \left\{ s\in \mathbb{V}_\epsilon:\prod_{i=1}^n |s_i| = g(\hat{p})\right\}.\] We observe that at any point of this level set,
   \begin{equation}\label{9/18/22.1} \frac{\partial g}{\partial s_j} = {\rm sgn}(s_j) \prod_{i\ne j} |s_i| = \frac{g(\hat{p})}{s_j}.\end{equation} The tangent space of $\mathcal{G}$ at $\hat{p}$ has normal vectors proportional to the gradient of $g(s)$ at $\hat{p}$ and, with respect to the $\{v_i\}_{i=1}^n$-basis, (\ref{9/18/22.1}) implies \begin{equation}\label{9/15/20.1} \nabla g(s)\Big|_{s=\hat{p}} = \,\, g(\hat{p}) \left(\frac{1}{p_1},\dots, \frac{1}{p_n}\right),\end{equation} which according to (\ref{3/18/12.1}) is proportional to $\nabla h(s)\big|_{s=\hat{p}}$.
Hence, $\mathcal{H}$ and $\mathcal{G}$ both contain $\hat{p}$, and share the same tangent space at $\hat{p}$, i.e., $\mathcal{G}$ and $\mathcal{H}$ are tangent at $\hat{p}$.  But we have seen that $\mathcal{H}$ is tangent to $\partial\Omega$ at $\hat{p}$.  Hence $\mathcal{G}$ is tangent to $\partial\Omega$ at $\hat{p}$, and $\hat{p}$ is thereby a critical point of $g$ on $\partial\Omega$.  The $n$-th root of $g$, which is the geometric mean, then also has a critical point on $\partial\Omega$ at $\hat{p}$, which is equivalent to the final statement in the theorem. 
 \end{proof} 
 
 By the way, the deterministic method is convergent when extended to the infinite dimensional setting, as shown in \cite{fgreensite:2022d} (that reference also provides an illustrative numerical example comparing it with standard zero-order Tikhonov regularization). 
 
 \bigskip
 
Theorem \ref{7/6/21.6} indicates that the regularized solution estimate is a critical point of the geometric mean on a particular discrepancy manifold.  With zero-order Tikhonov regularization, the regularized solution estimate is a critical point of the Euclidean norm on a discrepancy manifold (actually, it is the only relative minimum of the Euclidean norm on the discrepancy manifold) - and in generalized Tikhonov regularization, the minimization is at least performed with seminorms.  So we are already motivated to think of the geometric mean as some kind of a ``norm", albeit one not satisfying the triangle inequality. 

The geometric mean pertains to this regularization method precisely because of (\ref{9/18/22.1}), (\ref{9/15/20.1}), i.e., $\nabla g(s) = s^{-1}g(s)$ where $g(s)$ is the geometric mean raised to the $n$-th power (i.e., $g(s) = \prod_i |s_i|$), , and \begin{equation}\label{12/5/23.1} s^{-1}\equiv\left(\frac{1}{s_1},\dots, \frac{1}{s_n}\right).\end{equation} 
The geometric mean is a degree-1 positive homogeneous function extracted from $g(s) $, but we will modify it by multiplying it by a constant for reasons that will become clear shortly.   Thus, for a point $s = (s_1, \dots, s_n)\in \mathbb{R}^n$, with $s_i\ne 0$ for $i=1,\dots,n$, we define  \begin{equation}\label{12/6/23.3} \ell(s)\equiv\sqrt{n}\left(\prod_i |s_i|\right)^{1/n}.\end{equation}   Combining this with (\ref{12/5/23.1}), we have \begin{equation}\label{12/6/23.1}\ell(s)\ell(s^{-1}) = n = \mathbf{1}\cdot\mathbf{1}\equiv \|\mathbf{1}\|^2,\end{equation} where $\mathbf{1}\equiv (1,1,\dots,1)$.  Thus, for $s$ in some neighborhood of $\mathbf{1}$,
\[ \nabla \ell(s) = \frac{\ell(s)}{n}\left(\frac{1}{s_1},\dots,\frac{1}{s_n}\right) = \frac{s^{-1}}{\ell(s^{-1})},  \] i.e.,  \begin{equation}\label{12/5/23.2} s^{-1} = \ell(s^{-1})\nabla\ell(s),\end{equation}  Applying $\ell$ to both sides of the above equation, and using the degree-1 positive homogeneity of $\ell$,  we obtain \begin{equation}\label{12/6/23.5} \ell(\nabla\ell(s))
= 1.\end{equation}


 
 Equations (\ref{12/5/23.2}), (\ref{12/6/23.5}), describing the geometric mean, look a lot like the equations describing the Euclidean norm,  \begin{equation}\tag{\ref{11/29/23.5}}  s = \ell(s)\nabla\ell(s),\end{equation} and \begin{equation}\tag{\ref{11/29/23.6}} \ell(\nabla\ell(s)) = 1.\end{equation}  That is, the nonnegative solution for $\ell(s)$ from the latter two equations is the Euclidean norm, and the two equations express a point as the product of its Euclidean norm with an associated unit-norm direction.
So from this standpoint also, it seems that the geometric mean {\it is} indeed rather ``norm-like".  In particular, for the case of $\ell$ deriving from the geometric mean as in (\ref{12/6/23.3}) (which, in the above regularization context, we have already come to think of as some kind of a ``norm"), according to (\ref{12/5/23.2}), (\ref{12/6/23.5}) a point in a neighborhood of $\mathbf{1}$ has the decomposition,
\begin{equation}\label{8/4/20.2} s = \ell(s)\nabla \ell(s^{-1}) \mbox{  with   }  \ell(\nabla \ell(s^{-1})) = 1,\end{equation} where the expression $\nabla \ell(s^{-1})$ means that the gradient is obtained following which it is evaluated at $s^{-1}$.  Like the pair of equations (\ref{11/29/23.5}), (\ref{11/29/23.6}) for the Euclidean norm, the equations (\ref{8/4/20.2}) can also be interpreted as expressing a point as the product of its ``norm" with an associated unit-``norm"  direction.


Looking at (\ref{12/6/23.1}), (\ref{12/5/23.2}),  (\ref{12/6/23.5}) in isolation, we have to wonder: what is the significance of the ``inverse of a point", $s^{-1}$?  Well, it suggests that $s$ is an element of a unital algebra.  Which algebra?  In the case at hand, the inverse of a point happens to be given by (\ref{12/5/23.1}).   Equation (\ref{12/5/23.1}) indicates that the multiplicative identity element is ${\bf 1}= (1,1,\dots, 1)$, since this element will be equal to its inverse.  But we must have  $ss^{-1} = (s_1,\dots,s_n)(1/s_1,\dots,1/s_n) = {\bf 1}$, where element juxtaposition denotes the product. The equations of the last two sentences indicate that we must then be dealing with the algebra (whose vector space of elements is $\mathbb{R}^n$) defined by component-wise addition and multiplication, i.e., $\bigoplus_{i=1}^n \mathbb{R}$.    And that gets one thinking: how many other algebras are associated with an entity $\ell(s)$ satisfying something like (\ref{12/6/23.1}), (\ref{12/5/23.2}),  (\ref{12/6/23.5}), i.e., satisfying ``norm-like" equations respecting the multiplicative inversion operation of the algebra?
 
The answer is not many.  This is because, given (\ref{12/6/23.1}) (which mandates respect for the multiplicative inversion operation), the right-hand-side of (\ref{12/5/23.2}) is a gradient, i.e., $\nabla\log \ell(s)$ - but the left-hand-side, $s^{-1}$,  is not a gradient for most unital algebras.  
To rectify that, we would have to replace $s^{-1}$ on the left-hand-side of (\ref{12/5/23.2}) with $Ls^{-1}$, where $L$ has the feature that $Ls^{-1}$ {\it is} a gradient (and by the way, the left-hand-side of (\ref{12/6/23.5}) would then become $\ell(L^{-1}\nabla\ell(s))$ which is comparable to (\ref{11/29/23.2}), assuming $L$ to be nonsingular).   
\enlargethispage{1.5cm}

The requirement that $Ls^{-1}$ be a gradient means that $L$ must satisfy, \[ d([Ls^{-1}]\cdot ds) = 0.\]  We are thus led to Definition \ref{7/21/23.1} with $j=-1$, the beginning of our program.

\newpage



\section{Diagram of this paper}

\vspace{.5cm}
 
 \centerline{ \xymatrix{ & 
  *+[F--]{\txt{a new proof of the\\ Pythagorean Theorem\\ (also see arXiv:2209.14119)}} \ar@{-->}[d]^{motivation}_{alternative} & \txt{epimorphism problem} & \txt{ nonunital algebras} \\
  *+[F--]{\txt{an unusual \\ Inverse Problems\\ regularization \\ method (also see \\  arXiv:2209.14137)}} \ar@{-->}@/^1.3pc/[dr]_(.55){motivation\,} \ar@{-->}[r] ^(.45){motivation}& {\txt{Uncurling Metric Group \\ of type $j=-1$}} \ar[d]_(.45){maps}^(.43){to} \ar[dr]^(.59){part\,\,of}   \ar[r]^(.49){generalizes\,\,to}_(.45){(types\, j\ne 0,1\, are\, =\,?)} & *+[F]{\txt{Uncurling Metric Group \\ of type $j$, for all $j\in\mathbb{Z}$}}  \ar[ul]_{j=1} \ar[u]^{functor} \ar[ur]^{j >1} \ar[r]^<(.35){various}_<(.35){dimensions} & \txt{novel invariants} \ar@/^3pc/[ddl]   \\
  \txt{usual algebra norm} \ar@/^1.3pc/[r]_<(.45){generalizes\,\, to}   & *+[F]{\txt{Unital Norm Group}} \ar@/^1.5pc/[l]_(0.5){special \,unital\, norms}  \ar@/^1.2pc/[dl]_<(.45){cosmological\,\,\,\,}^<(.35){constant\, \ne\, 0} \ar[d]_(.45){triple}^(.45){index}   \ar[r]^<(.25){part\,\, of} & \txt{inversion characterization group}    \ar[d]_(.45){additional\,\, novel}^(.45){invariants?}  \\
  \txt{nonphysical cosmology\\ (not in this paper; see  \\   arXiv:2209.14127)} & \txt{novel invariants} \ar[r]  &  *+[F-:<15pt>]{\txt{coarse algebra classification}} \\
 & &  *+[F-,]{\txt{Algebra Classification Problem (wild)}} \ar[u]_{e.g.,\,Topology}^{motivation:} 
}
}

\vspace{1cm}
  
   \begin{thebibliography}{99}
   
 
   
    \bibitem{drozd:1972} Drozd YA (1972) Representations of commutative algebras. Functional Analysis and its Applications, 6:286-288.
   
       \bibitem{prest:1991} Prest M (1991) Wild representation type and undecidability.  Communications in Algebra, 19(3):919-929.
       
             \bibitem{sergeichuk:2000} Sergeichuk VV (2000) Canonical matrices for linear matrix problems.  Linear Algebra and its Applications, 317:53-102.

       
       \bibitem{belitskii:2005} Belitskii G, Lipyanski R, Sergeichuk VV (2005) Problems of classifying associative or Lie algebras and triples of symmetric or skew-symmetric matrices are wild.  Linear Algebra and its Applications, 407:249-262.
                            
  \bibitem{bourbaki:1989} Bourbaki N (1989) Algebra I, Springer, Berlin, p. 543.
  
              \bibitem{jacobson:1963} Jacobson N  (1963) Generic norm of an algebra, Osaka mathematics Journal, 15(1):25-50.
  
        \bibitem{schafer:1963}  Schafer R (1963) On forms of degree $n$ permitting composition, Journal of Mathematics and Mechanics, 12:777-792. 
         
      
       \bibitem{mccrimmon:2004} McCrimmon K (2004) A taste of Jordan algebras, Springer, New York.
       
 
       
              \bibitem{fgreensite:2022b} Greensite F (2022), A new proof of the Pythagorean Theorem and generalization of the usual algebra norm. arXiv:2209.14119
              
                    \bibitem{fgreensite:2022d} Greensite F (2022), Origin of the uncurling metric formalism from an Inverse Problems regularization method. arXiv:2209.14137
 
   
         \bibitem{birkhoff:1932} Birkhoff, GD (1932), A Set of Postulates for Plane Geometry (Based on Scale and Protractors), Annals of Mathematics, 33: 329–345.
         
               \bibitem{vogel:2002} Vogel C (2002) Computational Methods for Inverse Problems. Philadelphia: SIAM.
      
      \bibitem{engl:2000}  Engl HW,  Hanke M,  Neubauer A (2000) Regularization of Inverse Problems, Klewer, Boston.
      
      \enlargethispage{1cm}
      
      \bibitem{hansen:1992} Hansen PC (1992) Analysis of discrete ill-posed problems by means of the L-curve, SIAM Review, 34:563-580.
      
       
      \end{thebibliography}
  
  
   \end{document}   
 